\section{Depth reduction for interacting walks}\label{sec:depth}

\Cref{thm:D} reduces the oblivious case of~(T) to two hypotheses. This section is about the structural one: deleting a
small fraction of the edges of the walk DAG of a head schedule so that every remaining path contains few long hops.
\Cref{thm:B} does this for home-based walks, for every coupling. Here we state the general problem at a fixed
fraction~$\eps$ (\cref{prob:tcp}), rule out packing certificates for lower bounds (\cref{prop:packing}), give
evidence against per-head band bookkeeping (\cref{rem:dep-N,rem:dep-Nprime}, proof sketches only), and study what we
regard as the sharpest open instance: two deeply nested walks, shifted against each other in time (\cref{def:dnp}). One
such walk alone is easy (\cref{lem:dep-visit}), and the pair with $m$ levels of nesting is depth-reducible: clock labels cut
it at a depth exponential in~$m$ (\cref{prop:dep-clock}). On a cyclic cover of the pair the cost of every label obeys exact
identities (Lemma~Zq, \cref{lem:chain,lem:transfer}); with them we prove, for the pair with one level of nesting, that
the labels that count exactly along one head pay a price per hop on the cover that does not depend on the depth
(\cref{thm:E,thm:E1}). Whether the pairs are depth-robust as $m$ grows remains open (\cref{conj:Aprime}).

\subsection{The problem}\label{sec:dep-problem}

We use the walk DAG of \cref{sec:dec-B}: for walks $x_1,\dots,x_k$ of length~$T$, its vertices are the times
$[0,T)$, its edges are the state edges $t\to t-1$ and the long hops $u\to\prev_h(u)$, a pause is identified with the
state edge joining the same two times, and an edge is determined by its endpoints. A \emph{cut} is a set of edges; it
has \emph{depth at most~$D$} if every path that uses no edge of the cut contains at most $D$ long hops.

\begin{problem}[$\TCP(k)$ at fixed $\eps$]\label{prob:tcp}
Let $k\ge1$. Is there, for every $\eps>0$, an integer $D=D(k,\eps)$ such that for every~$T$ and all walks
$x_1,\dots,x_k$ of length~$T$, coupled arbitrarily, the walk DAG has a cut with at most $\eps T$ edges and depth at
most~$D$?
\end{problem}

A family of $k$-tuples of walks is \emph{depth-robust} if for some $\eps>0$ no single~$D$ serves all its members.
Cuts are conveniently described by labels. For an integer $D\ge0$ and a map $\Psi\colon[0,T)\to\{0,\dots,D\}$, call a
long hop $u\to p$ \emph{violated} if $\Psi(u)\le\Psi(p)$, call a time $t\ge1$ a \emph{descent} if
$\Psi(t)<\Psi(t-1)$, and let $\cost(\Psi)$ be the number of violated long hops plus the number of descents.

\begin{lemma}[Cuts are labels]\label{lem:dep-label}
\notformalized
For all walks $x_1,\dots,x_k$ of length~$T$ and every $D\ge0$, the smallest size of a cut of depth at most~$D$ equals
the smallest value of $\cost(\Psi)$ over all maps $\Psi\colon[0,T)\to\{0,\dots,D\}$.
\end{lemma}

\begin{informalproof}
Given $\Psi$, let $F$ consist of the violated long hops and of the state edges $t\to t-1$ at descents~$t$; so
$\card F=\cost(\Psi)$. Along every edge $a\to b$ outside~$F$ we have $\Psi(b)\le\Psi(a)$, with strict inequality for
long hops. Along a path avoiding~$F$, $\Psi$ therefore never increases and drops at every long hop; as it takes values
in $\{0,\dots,D\}$, the path contains at most $D$ long hops. Conversely, let $F$ be a cut of depth at most~$D$, and
let $\Psi(t)$ be the largest number of long hops on a path that starts at~$t$ and avoids~$F$; so
$0\le\Psi(t)\le D$. For an edge $a\to b$ not in~$F$, prepending it to a path from~$b$ gives
$\Psi(a)\ge\Psi(b)+1$ if the edge is a long hop and $\Psi(a)\ge\Psi(b)$ otherwise. Hence every violated long hop and
every state edge at a descent of~$\Psi$ belongs to~$F$, and $\cost(\Psi)\le\card F$.
\end{informalproof}

\paragraph{Why the problem matters.} \Cref{prob:tcp} is the walk-level form of the structural hypothesis of \cref{thm:D}. If it
has a positive answer for~$k$, the proof of \cref{cor:B}, with the cut provided by \cref{prob:tcp} in place of the one
provided by \cref{thm:B}, shows that $\mathrm{depth}\text{-}\TCP(k,C,\beta,\eps,D)$ holds for all $C\ge1$,
$\beta\ge2$ and $0<\eps\le1$, with $D=(D(k,\eps/(4\beta C))+1)\lceil4\beta C/\eps\rceil$; \cref{thm:D} then reduces
the oblivious case
of~(T) to the crux alone, as \cref{cor:B} does for home-based heads. The problem is posed for all walks, and not only
for the head schedules of transposers, because a machine can run any schedule that its finite control generates on
extra work tapes without changing its output; such dummy traffic enters the dependency DAG. A family that
oblivious machines can generate in linear time, and that is depth-robust at every $\eps$ below the crux rates in
question, would therefore make the structural hypothesis fail for some machines, and the route of \cref{thm:D} would
need a version of it that sees only the part of the DAG able to carry the transposition. On the
construction side, \cref{thm:D} says that an oblivious linear-time transposer either beats the crux or has a head
schedule without genies of small cost and bounded depth; up to constants this means that its walk DAG has no cut of
size $\eps T$ and bounded depth, since a genie~$\Gamma$ yields a cut of at most $(k+1)\card\Gamma$ edges and the
conversions of \cref{sec:dec-corB} go the other way. The walks of this section show what the second alternative could
look like.

For $k=1$ the problem is easy. Call a long hop $(p,u)$ of a single walk~$x$ an upper hop if $x>x(p)$ on $(p,u)$ and a
lower hop if $x<x(p)$ there. Two upper hops never cross as intervals of time: if $p<p'<u<u'$, then $x(p')>x(p)$
because $p'\in(p,u)$, while $u\in(p',u')$ and $x(u)=x(p)<x(p')$, contradicting that $x>x(p')$ on $(p',u')$; the same
holds for lower hops. Every long hop is an upper or a lower hop, since $x\ne x(p)$ on $(p,u)$ and $x$ moves by at most~$1$ per
step. Drawing the upper hops on one side of the time axis and the lower hops on the other shows that
the walk DAG of one walk is planar; its maximum degree is at most~$4$. By the planar separator theorem of Lipton and
Tarjan~\cite{LiptonTarjan1979}, applied recursively as in~\cite{LiptonTarjan1980}, some set of $O(T/\sqrt M)$ vertices
leaves components of at most $M$ vertices. Cutting the at most four edges at each of these vertices cuts
$O(T/\sqrt M)$ edges, and a path that avoids the cut stays inside one component, so it has at most $M-1$ long hops;
this gives $D(1,\eps)=O(\eps^{-2})$. For $k\ge2$ the union of the walks' hop families is in general not planar, and only special
classes are solved, chief among them the home-based walks of \cref{thm:B}.

\paragraph{Relation to segregators.} Paul, Pippenger, Szemer\'edi and Trotter~\cite{PPST1983} showed that the
computation graphs of multitape machines have small segregators: deleting $O(n/\log^*n)$ vertices leaves every vertex
with at most $O(n/\log^*n)$ predecessors. This is the combinatorial core of their separation of deterministic from
nondeterministic linear time. \Cref{prob:tcp} is a depth version of the same problem with constant parameters.
Through the conversions of \cref{sec:dec-corB}, a positive answer gives genies of cost $\eps L$ and bounded
depth~$D'$, and by the cone lemma (\cref{lem:cone}) every step then has at most $(k+1)^{D'+1}$ predecessors reached
without passing through the genie: the deleted set is a fixed fraction of the steps, far more than $n/\log^*n$, but
the remaining cones have bounded size instead of size $n/\log^*n$. Santhanam stated that if, for every~$r$, the
multi-pushdown graphs $H_r(n)$, a class containing the computation graphs of multitape machines, had segregators
deleting $o(n/\log n)$ vertices and leaving every vertex at most $o(n)$ predecessors, then deterministic and
nondeterministic time $n\log n$ would differ~\cite[Thm.~2.7]{Santhanam2001}, and indicated that its proof is analogous
to that of his Theorem~2.5. Since the deleted
part in \cref{prob:tcp} is linear, a positive answer would not by itself give a new separation of complexity classes.

\subsection{Uniform recursive walks}\label{sec:dep-walks}

The instances of this section are built from one family of walks in which every level of a recursive nesting
receives the same total time. Since \cref{prob:tcp} concerns all schedules, these walks are relevant even though the
results of \cref{sec:motion} constrain the macroscopic motion of transposers.

\begin{definition}[Uniform recursive walks]\label{def:dep-walk}
Let $g\ge1$ and $m\ge0$ be integers, and let $s_0$ be an integer with $s_i:=s_0-ig\ge1$ for $0\le i\le m$. A
\emph{level-$m$ stroke} of direction $\sigma\in\{-1,1\}$ is the sequence of $2^{s_m}$ steps equal to~$\sigma$. For
$i<m$, a \emph{level-$i$ stroke} of direction~$\sigma$ is the concatenation, for $j=0,\dots,2^{g-1}-1$, of a
\emph{segment} of $2^{s_{i+1}+1}$ steps equal to~$\sigma$ followed by a \emph{V}: a level-$(i+1)$ stroke of
direction~$-\sigma$ (the \emph{back stroke}) followed by a level-$(i+1)$ stroke of direction~$\sigma$ (the
\emph{forward stroke}). The walk $U(s_0,m,g)$ is the walk $x\colon\Z\to\Z$ with $x(0)=0$, periodic with period
$P:=4(m+1)2^{s_0}$, whose steps $x(t+1)-x(t)$, $t\in[0,P)$, are those of four level-$0$ strokes of directions
$1,-1,-1,1$, in this order.
\end{definition}

By induction on $m-i$, a level-$i$ stroke has $(m-i+1)2^{s_i}$ steps and moves the head by $\sigma2^{s_i}$: its
segments contribute $2^{g-1}\cdot2^{s_{i+1}+1}=2^{s_i}$, and each V returns the head to where it started. So every
period has $4\cdot2^{s_0}=P/(m+1)$ steps belonging to segments of level-$i$ strokes for each $i<m$, and as many
belonging to level-$m$ strokes, and the walk never pauses. Each step belongs to exactly one segment or level-$m$
stroke; we say that the stroke emits it, and a time~$t$ belongs to the stroke that emits the step from $t$ to $t+1$.
The \emph{span} of a stroke or of a V is the set of times belonging to it, and its \emph{end time} is the first time
after the span. The terms \emph{during}, \emph{first time} and \emph{last time} refer to spans.

If a level-$i$ stroke of direction~$\sigma$, $i<m$, starts at cell~$a$, its segment~$j$ ends at the \emph{grid point}
$b_j=a+\sigma(j+1)2^{s_{i+1}+1}$, and the V that follows is \emph{attached} at~$b_j$. By
\cref{lem:dep-structure}(a) below, this V stays in its \emph{range}, the interval of cells between its \emph{base}
$b_j$ and its \emph{valley cell} $b_j-\sigma2^{s_{i+1}}$; the \emph{level} of the V is the level of its strokes, its
\emph{apex time} is the first time of its back stroke and its \emph{valley time} the first time of its forward
stroke. The \emph{open half-segments} of the stroke are the open intervals of cells between consecutive points of
$a+\sigma2^{s_{i+1}}\Z$ that lie between $a$ and $a+\sigma2^{s_i}$; the \emph{open range} of a level-$i$ stroke is the set of
cells strictly between $a$ and $a+\sigma2^{s_i}$, and for a level-$m$ stroke we use it in place of the open
half-segments; the open range of a V is its range without the base and the valley cell. A \emph{turn} is a time~$t$ with
$x(t+1)-x(t)\ne x(t)-x(t-1)$, and its \emph{level} is the level of the stroke to which $t$ belongs. Intervals of
cells are written without regard to orientation, and ``beyond~$b$ in direction~$\sigma$'' means $\sigma(c-b)>0$.

\begin{lemma}[Structure of $U$]\label{lem:dep-structure}
\notformalized
Let $x=U(s_0,m,g)$.
\begin{enumerate}[label=(\alph*)]
\item During a level-$i$ stroke~$S$ of direction~$\sigma$ starting at cell~$a$, including its end time, the head
stays between $a$ and $a+\sigma2^{s_i}$, and it is at $a+\sigma2^{s_i}$ at the end time. The first and the last step
of~$S$ equal~$\sigma$. If $i<m$, the part of~$S$ up to the end time of the V at~$b_j$ stays between $a$ and~$b_j$, and
the part after it stays between $b_j$ and $a+\sigma2^{s_i}$. During a V, including its end time, the head stays in
the range of the V.
\item The turns are exactly the apex and valley times of all V's, each of the level of its V, and, in every period,
the first times of the second and of the fourth level-$0$ stroke, which have level~$0$; at these two times the head
is at $2^{s_0}$ and at $-2^{s_0}$.
\item During a V of level~$i$, every cell strictly inside its range is visited exactly $2(m-i+1)$ times and each of
its two endpoints exactly $m-i+1$ times. In every period, every cell~$c$ with $\abs c<2^{s_0}$ is visited exactly
$2(m+1)$ times and each of $\pm2^{s_0}$ exactly $m+1$ times.
\item For every stroke~$S$ whose first time~$t_S$ is nonnegative, the number of times $t\in[0,t_S)$ with $x(t)=c$
is the same for all cells~$c$ of any one open half-segment of~$S$ (of the open range of~$S$ if $S$ has level~$m$).
\item Let $1\le j\le m$. If no turn of level less than~$j$ lies strictly inside a hop $(p,u)$, then
$\abs{x(t)-x(p)}\le2^{s_j}$ for all $t\in[p,u]$.
\end{enumerate}
\end{lemma}

\begin{informalproof}
(a) Induction on $m-i$. A level-$m$ stroke is monotone. For $i<m$, the segments move the head monotonically from~$a$
to $a+\sigma2^{s_i}$, reaching $b_j$ at the end of segment~$j$. By the induction hypothesis, the back stroke of the V
at~$b_j$ stays between $b_j$ and $b_j-\sigma2^{s_{i+1}}$ and ends at the valley cell, and the forward stroke stays in
the same range and ends at~$b_j$. This range lies between $b_j-\sigma2^{s_{i+1}+1}$, the start of segment~$j$, and
$b_j$, which gives all the claims about~$S$. The first step of~$S$ is a segment step, and its last step is the last
step of its last forward stroke, which equals~$\sigma$ by induction.

(b) Inside a segment and inside a level-$m$ stroke all steps are equal. Let $t$ be a time at which the steps from
$t-1$ to~$t$ and from $t$ to $t+1$ are emitted by different strokes. Let $Y$ be the largest segment or stroke whose
span ends at $t-1$, and $Z$ the largest segment or stroke whose span begins at~$t$. They are consecutive parts of one
stroke or of the sequence of level-$0$ strokes, so the pair $(Y,Z)$ is of one of four kinds, and by~(a) the last step
of~$Y$ and the first step of~$Z$ are their directions. (1)~$Y$ is a segment of a stroke of direction~$\sigma$ and $Z$
the back stroke of the V attached at its end: the steps are $\sigma$ and~$-\sigma$, and $t$ is the apex time, a turn
of the level of the back stroke. (2)~$Y$ is a back stroke and $Z$ the forward stroke of the same V: the steps are
$-\sigma$ and~$\sigma$, and $t$ is the valley time, a turn of the level of the forward stroke. (3)~$Y$ is a forward
stroke and $Z$ the next segment of the stroke to which its V is attached: both steps are~$\sigma$, and $t$ is not a
turn. (4)~$Y$ and~$Z$ are consecutive level-$0$ strokes: their directions $1,-1,-1,1$, followed by~$1$ in the next
period, give a turn exactly when $Z$ is the second or the fourth stroke of a period, and then the head is at
$2^{s_0}$ or $-2^{s_0}$ by~(a).

(c) By induction on $m-i$ we prove the claim about V's. A V of level~$m$ consists of two monotone runs of $2^{s_m}$
steps in opposite directions: the back stroke visits every cell of the range except the valley cell, the forward
stroke every cell except the base. Let $i<m$, and suppose by symmetry that the back stroke goes down, so that the
range is $[v,b]$ with $v=b-2^{s_i}$. Put $s=s_{i+1}$ and $K=2^{g-1}$, and let $v=\gamma_0<\gamma_1<\dots<\gamma_K=b$
be the points of $v+2^{s+1}\Z$ in $[v,b]$; since $2^{s_i}=K2^{s+1}$, these are also the points of $b-2^{s+1}\Z$ in
$[v,b]$. The segments of the back stroke visit $b,b-1,\dots,v+1$ once each, and its V's have level $i+1$, are attached
at $\gamma_{K-1},\dots,\gamma_0$ and have ranges $[\gamma_k,\gamma_k+2^s]$. The segments of the forward stroke visit
$v,v+1,\dots,b-1$ once each, and its V's are attached at $\gamma_1,\dots,\gamma_K$ and have ranges
$[\gamma_k-2^s,\gamma_k]$. By the induction hypothesis each of these V's visits the cells strictly inside its range
$2(m-i)$ times and its endpoints $m-i$ times, and the spans of the segments and of these V's partition the span of the
V. A cell strictly between $\gamma_k$ and $\gamma_k+2^s$, or strictly between $\gamma_k+2^s$ and $\gamma_{k+1}$, lies
strictly inside exactly one of the attached V's and is visited $1+1+2(m-i)$ times. The cell $\gamma_k+2^s$
($0\le k<K$) is an endpoint of two attached V's, the back stroke's at $\gamma_k$ and the forward stroke's at
$\gamma_{k+1}$, and so is every $\gamma_k$ with $0<k<K$; these cells are visited $1+1+2(m-i)$ times. The base~$b$ is
visited once by the back stroke's segments and $m-i$ times by the forward stroke's V at $\gamma_K=b$, and the valley
cell~$v$ once by the forward stroke's segments and $m-i$ times by the back stroke's V at $\gamma_0=v$. This proves
the claim for V's. In a period, the first two level-$0$ strokes have, over $[0,2^{s_0}]$, the same segments and
attached V's as the forward and back strokes of a V over this range (in the other order), so the count for V's applies
verbatim and shows that they visit every $c$ with $0<c<2^{s_0}$ exactly $2(m+1)$ times and $0$ and $2^{s_0}$ exactly $m+1$
times; the last two strokes do the same over $[-2^{s_0},0]$.

(d) For a stroke~$S$ with $t_S\ge0$ write $\mathrm{pre}_S(c)$ for the number of times $t\in[0,t_S)$ with $x(t)=c$.
By~(a) and~(c), during a level-$i$ stroke with $i<m$ the number of visits to~$c$ is constant on each open
half-segment: the segments visit each cell between the first cell (included) and the last cell (excluded) once, and
the V at~$b_j$ adds $2(m-i)$ visits on the open half-segment of segment~$j$ next to~$b_j$ and none elsewhere inside
open half-segments. During a level-$m$ stroke every cell of the open range is visited once. We now argue down the
tree of strokes; since time~$0$ starts a level-$0$ stroke, the parent of every stroke~$S$ with $t_S\ge0$ also starts at a
nonnegative time. Let $S$ be a level-$0$ stroke of period~$n$. If it is the first or the third stroke of the period,
$\mathrm{pre}_S(c)=2(m+1)n$ on its open range by~(c). If it is the second (fourth) stroke, $\mathrm{pre}_S$ is the sum
of $\mathrm{pre}_{S'}$ and of the visits during~$S'$, where $S'$ is the first (third) stroke of the period, and $S$
and~$S'$ have the same open half-segments because $2^{s_0}$ is a multiple of $2^{s_1}$ (for $m=0$, the same open range). Next let $S$ be the back
stroke of a V attached at~$b_j$ to a stroke~$S'$. The open range of the V is an open half-segment of~$S'$, and on it
$\mathrm{pre}_S=\mathrm{pre}_{S'}+1$: by~(a) the part of~$S'$ before segment~$j$ stays between the first cell of~$S'$
and the start of segment~$j$, and segment~$j$ visits each cell of the open range once. So $\mathrm{pre}_S$ is constant
on the open range of~$S$, which contains all its open half-segments. If $S$ is the forward stroke of that V and $B$ its
back stroke, $\mathrm{pre}_S$ is the sum of $\mathrm{pre}_B$ and of the visits during~$B$, and $B$ and~$S$ have the
same open half-segments, or the same open range at level~$m$, because their common displacement is a multiple of the
spacing of their half-segments. Hence $\mathrm{pre}_S$ is constant on the open half-segments, or the open range, of every
stroke.

(e) Let $c=x(p)$. As $u-p\ge2$ and the walk moves by $\pm1$, the first time $\tau^*\in(p,u)$ at which $\abs{x(t)-c}$
is largest is a turn, and it suffices to show $\abs{x(\tau^*)-c}\le2^{s_j}$. The level of~$\tau^*$ is at least
$j\ge1$, so by~(b) $\tau^*$ is the apex or valley time of a V of level at least~$j$, and therefore belongs to the span
of a V of level exactly~$j$, say~$\mathcal V$: a V of level greater than~$j$ lies inside a stroke of level at
least~$j$, and such a stroke lies inside a V of level~$j$. Let $\mathcal V$ be attached to the stroke~$S'$, of
level~$j-1$ and direction~$\sigma$, at its base~$b$, so that its range is $[v,b]$ with $v=b-\sigma2^{s_j}$. By~(a),
$x(\tau^*)\in[v,b]$, so it suffices to show $c\in[v,b]$. Suppose not. By~(a) the head is in $[v,b]$ during
$\mathcal V$ and at its end time~$t_E$, so $p$ precedes the apex time~$t_A$ of~$\mathcal V$ and $u$ follows~$t_E$.

\emph{Case 1: $c$ lies beyond~$b$ in direction~$\sigma$.} By~(a) the part of~$S'$ before~$\mathcal V$ stays between
the first cell of~$S'$ and~$b$, so $p$ precedes the first time of~$S'$. If $j\ge2$, $S'$ is a back or forward stroke
of a V of level $j-1$, whose first time is an apex or valley time, a turn of level $j-1$ in $(p,t_A)$. If $j=1$ and
$S'$ is the second or fourth level-$0$ stroke of its period, its first time is a turn of level~$0$ in $(p,t_A)$. If
$S'$ is the first (third) level-$0$ stroke, then $\sigma=1$ and $c>b>0$ ($\sigma=-1$ and $c<b<0$); the two preceding
level-$0$ strokes stay in $[-2^{s_0},0]$ (in $[0,2^{s_0}]$), so $p$ precedes both of them, and the first time of the
later one, the fourth (second) level-$0$ stroke, is a turn of level~$0$ in $(p,t_A)$.

\emph{Case 2: $c$ lies beyond~$v$ in direction~$-\sigma$.} Let $A_{j-1}=S'$, and for $1\le l\le j-1$ let $A_{l-1}$ be
the stroke to which the V containing $A_l$ is attached; $A_l$ has level~$l$. From~$t_E$ on, by~(a), the head stays
beyond~$b$ in direction~$\sigma$, or at~$b$, until the end time of~$A_{j-1}$. If $A_{j-1}$ is a back stroke, its end
time is the valley time of its V, a turn of level $j-1$. If $A_{j-1}$ is a forward stroke, it has the direction
$\sigma$ of $A_{j-2}$, its end time is the end time of a V attached to~$A_{j-2}$ at the last cell of~$A_{j-1}$, and
by~(a) the head stays beyond that cell, hence beyond~$b$, until the end time of~$A_{j-2}$. Repeating this, if some
$A_l$ with $1\le l\le j-1$ is a back stroke and $l^*$ is the largest such~$l$, the head stays beyond~$b$ (or at~$b$)
from~$t_E$ until the end time of~$A_{l^*}$, which is a valley time, a turn of level $l^*\le j-1$. Otherwise all
$A_l$ have direction~$\sigma$, and the head stays beyond~$b$ (or at~$b$) until the end time of the level-$0$
stroke~$A_0$. If $A_0$ is the first (third) stroke of its period, that end time is a turn of level~$0$. If it is the
second (fourth) stroke, the next level-$0$ stroke continues in direction~$\sigma$ and stays beyond~$b$, and its end
time is a turn of level~$0$. In every case the turn follows~$t_E$, and the head is not at~$c$ between $t_E$ and the
turn, so the turn precedes~$u$.

In both cases a turn of level less than~$j$ lies strictly inside $(p,u)$, contrary to the hypothesis.
\end{informalproof}

\subsection{Union labels and the nesting barrier}\label{sec:dep-nesting}

The labels of \cref{thm:B} are lexicographic: they carry counts of events at a few scales and compare them from the
top. A natural way to extend this to arbitrary walks is to let each head decide, at each cell and time, which bands of
scales (groups of consecutive scales; in \cref{sec:dec-B} a band was a single scale) it currently needs, and to carry, for every band needed by some head, a count of the events of that band (a
\emph{union label}). Two requirements make such a label work. A hop must find its own band among the tracked bands at
both of its ends, so that it is decided there; and a change of the bands tracked by one head breaks the comparison
along the hops of the other heads that span it, of which there can be about $2^\theta$ when $\theta$ is the bottom
scale of the changed band, so changes at high bands must be rare. We make the per-head part precise and sketch, without a
full proof, why it cannot be met with bounded load on deeply nested walks. We do not claim that such rules would
suffice for \cref{prob:tcp}; the point is that they appear not to exist.

\begin{definition}[Band rules]\label{def:dep-band}
Let $x$ be a walk of length~$T$ without pauses. For a turn~$\tau$ at which $x$ has a local maximum, a time~$t$ is
\emph{higher} than~$\tau$ if $x(t)>x(\tau)$, or $x(t)=x(\tau)$ and $t<\tau$; the \emph{prominence} $\mathrm{pr}(\tau)$
is the least $d\ge1$ such that on one side of~$\tau$ the walk reaches a time higher than~$\tau$ without going below
$x(\tau)-d$ on the way, and $x(\tau)-\min x$ if no time is higher than~$\tau$. For local minima exchange above and
below, an earlier time at the same height counting as lower. Fix integers $h_0,W\ge1$ and offsets $r\in[0,h_0)$,
$r'\in[0,W)$. A scale $\theta\ge0$ lies in \emph{band} $\lfloor(\lfloor(\theta+r)/h_0\rfloor+r')/W\rfloor$, and
$\theta(b)$ is the least scale of band~$b$. The band of a turn is the band of $\lfloor\log_2\mathrm{pr}(\tau)\rfloor$.
The \emph{owner} of a hop $(p,u)$ is the first time in $(p,u)$ at which $\abs{x(t)-x(p)}$ is largest, a turn, and the
\emph{own band} of the hop is the band of its owner. For a band $b\ge1$, the \emph{$b$-gaps} are the intervals
$[\tau,\tau')$ between consecutive turns of band at least~$b$ (and the interval before the first such turn). A
\emph{band rule} assigns to every cell~$c$ and time~$t$ a set $\mathcal M(c,t)$ of bands containing~$0$. It is
\emph{consistent} if for every $b\ge1$ the membership $b\in\mathcal M(c,t)$ depends only on~$c$ and on the $b$-gap
containing~$t$; its \emph{load} is $\max_t\card{\mathcal M(x(t),t)}$; $\mathrm{Viol}$ is the number of hops whose own
band is missing from $\mathcal M(x(p),p)$ or from $\mathcal M(x(u),u)$; and
\[
  D_{\mathrm{sum}}=\sum_{t}\;\sum_{b\ge1,\ b\in\mathcal M(x(t),t)\,\triangle\,\mathcal M(x(t+1),t+1)}2^{\theta(b)},
\]
while $D_{\max}$ counts, for each~$t$, only the largest band of the symmetric difference.
\end{definition}

A union label for $k$ heads tracks, at time~$t$, the union of the sets $\mathcal M_h(x_h(t),t)$ of the heads; with
load~$P_0$ per head it has at most $kP_0$ digits, so its depth is bounded in terms of $k$, $P_0$ and the moduli. Its
cost is bounded by the heads' violations, by their changes weighted as above (a change at band~$b$ is charged to the
hops of other heads that span it), and by wrap terms. This is the setting of the following two statements.

\begin{remark}[The nesting barrier for consistent band rules]\label{rem:dep-N}
Proof sketch only: we state the result precisely, but its proof exists only at the level of the following sketch;
exact linear programs support it numerically. Let
$g\ge\max(Wh_0,2)$ and $s_m\ge Wh_0+1$, let $x=U(s_0,m,g)$ on a whole number of periods, $T$ steps in all, and put
$M=T/(m+1)$. Then for every choice of the offsets, every consistent band rule of load at most~$P_0$ satisfies
\[
  \mathrm{Viol}+\gamma D_{\mathrm{sum}}\;\ge\;(m+2-P_0)M-2^{1-g}T-2^{1-s_m}M-O(m2^{s_0}),\qquad\gamma=2^{Wh_0-g};
\]
if the rule need not contain band~$0$, replace $m+2-P_0$ by $m+1-P_0$. \emph{Sketch.} Up to boundary effects the
prominences of the turns of~$U$ are $2^{s_0+1}$ and $2^{s_1},\dots,2^{s_m}$, one value per level of nesting (at a
double extremum the later turn carries the prominence of the deeper level), and for every offset their bands
$\beta_0>\beta_1>\dots>\beta_m$ are distinct, with $\theta(\beta_i)\ge s_i-Wh_0+1$. For a level-$i$ stroke~$S$ let
$\mathrm{gap}(S)$ be the $\beta_i$-gap containing its first segment step, and let $y(c,S)\in\{0,1\}$ be the membership
of~$\beta_i$ at $(c,\mathrm{gap}(S))$, well defined by consistency. (1)~For every cell~$c$ that the segments of~$S$
visit inside $\mathrm{gap}(S)$, the first visit of~$c$ in $\mathrm{gap}(S)$ ends a hop of own band~$\beta_i$, which
$y(c,S)=0$ violates; distinct pairs $(c,S)$ give distinct hops, except for a fraction at most $q_02^{-g}$ of each
level, $q_0=1/(2(1-2^{-g}))$, which comes from the last V's of strokes whose gap ends at their own arrival, and except
for the first cell of each back stroke. (2)~Every switch of $y(\cdot,S)$ between adjacent cells is a change of
$\beta_i$ at a segment step inside $\mathrm{gap}(S)$, of weight $2^{\theta(\beta_i)}\ge2^{1-Wh_0}\cdot2^{s_i}$, and
distinct pairs (stroke, switch) give distinct pairs (step, band). (3)~A cell~$c$ of a level-$m$ stroke with ancestors
$S_0\supset S_1\supset\dots\supset S_m$ is visited by the segments of~$S_m$ at a time that lies in $\mathrm{gap}(S_i)$
for every~$i$, except at exceptional levels, where $S_m$ lies in the last V of an ancestor whose gap ends early; so
the load gives $\sum_iy(c,S_i)\le P_0-1+\#\{\text{exceptional levels of }c\}$, and the expected number of exceptional
levels of a uniform such cell is at most $2q_02^{-g}m$. (4)~Averaging~(3) over a uniform level-$m$ stroke and a
uniform cell of it, which gives every block (a segment with its V) of every ancestor the same mass, and comparing
with~(1) and~(2) yields the bound; the exceptional terms add up to at most $2^{1-g}T$ because $g\ge2$.

Per head, a union label would need, heuristically, $\mathrm{Viol}\le\eps(N+T)$ and $D_{\max}\le\eps(N+T)$, where $N\le T$ is the number of
long hops. As $D_{\mathrm{sum}}\le D_{\max}/(1-2^{-Wh_0})$, the bound excludes this for all $P_0$, $W$ and~$h_0$ once
$\eps<0.37$ (take $g=Wh_0+4$, $m=10P_0$ and $s_m$ large), and once $\eps<0.21$ at $g=Wh_0+2$. Numerically, the exact
optimum of $\mathrm{Viol}+\gamma D_{\mathrm{sum}}$ over all consistent rules, computed by linear programming on
seventeen instances with $g\in\{1,2,3,4\}$ and $m\le5$, was $(m+2-P_0)M'-1$, where $M'$ is the mass per level in the
window used; the bound is loose only by its correction terms.
\end{remark}

\begin{remark}[The nesting barrier without consistency]\label{rem:dep-Nprime}
Proof sketch only: under the hypotheses of \cref{rem:dep-N}, let $t\mapsto\mathcal M(t)$ be any choice of sets of bands
containing~$0$, with $\card{\mathcal M(t)}\le P_0$ and no consistency requirement; it may depend on the whole walk, and
it may be randomised or fractional. Let $\mathrm{Viol}_u$ be the number of hops whose own band is missing from
$\mathcal M(u)$ at the later endpoint~$u$, and let $D_{\mathrm{sum}}$ count the changes of non-zero bands with weights
$2^{\theta(b)}$. Then $\mathrm{Viol}_u+\gamma D_{\mathrm{sum}}$ satisfies the bound of \cref{rem:dep-N}. Every band rule,
read along the walk as $\mathcal M(t)=\mathcal M(x(t),t)$, is such a choice, with the same load and changes and with
$\mathrm{Viol}_u\le\mathrm{Viol}$; so relaxing consistency does not help, and the thresholds $0.37$ and $0.21$ remain.
\emph{Sketch.} Steps~(2) and~(3) of the sketch of \cref{rem:dep-N} are the only uses of consistency, and both are replaced
by a trichotomy in time. For a level-$i$ stroke~$S$ and its block~$j$ (segment~$j$ with its V), let $J$ be the
block's time window, from just after the start of segment~$j$ to the start of segment~$j+1$ (the first and last
windows extend to the ends of $\mathrm{gap}(S)$); the windows of one level are disjoint. Exactly one of the following
holds: $\beta_i$ changes at a step inside~$J$, at cost $\gamma2^{\theta(\beta_i)}\ge2^{s_{i+1}+1}$, the number of
cells of segment~$j$; $\beta_i$ is absent throughout~$J$, so all the hops entering the cells of segment~$j$, whose
later endpoints lie in~$J$, are violated; or $\beta_i$ is present throughout~$J$, in particular at the deepest visits
inside the V of block~$j$, where it uses load. The load at the deepest visit of a cell then bounds the number of
levels~$i$ of the third kind by $P_0-1$, up to exceptional levels, and the averaging of step~(4) concludes. A
layer-cake decomposition of $1-\mathcal M$ gives a second proof, as a max-flow certificate. Numerically, when the own
band is required at both endpoints of a hop, the optimum over such choices of sets was lower than over consistent
rules only by a boundary term on the instances computed, and a bound with only boundary losses is conjectured; when it
is required at the later endpoint alone, the optimum was lower by a small multiple of~$M$.
\end{remark}

\paragraph{Why they matter.} \Cref{rem:dep-N,rem:dep-Nprime} indicate, with proof sketches only, that the union-label
route to \cref{prob:tcp} fails in its natural form: per-head band bookkeeping with load independent of the depth of
nesting, own-band decisions and $2^\theta$-weighted changes would fail already on a single walk, offline and for every
offset. The obstruction does not
concern $\TCP(1)$, which holds by planarity; what fails is a per-head format that would have handled all couplings at
once. A repair must decide some hops at a band other than their own, as the labels of \cref{thm:B} do, or weigh
changes differently, which is legitimate only against other heads that are not adversarial, or use a different
mechanism altogether.

\paragraph{One nested walk is nevertheless easy.} A single digit that ignores bands handles all levels of one walk at
once.

\begin{lemma}[One nested walk is easy]\label{lem:dep-visit}
\notformalized
Let $x=U(s_0,m,g)$ on $[0,T)$, where $T$ is a multiple of~$P$, let $q\ge1$, and let $\kappa(t)$ be the number of times
$s\in[0,t)$ with $x(s)=x(t)$. For $\omega\in\{0,\dots,q-1\}$ put $\Psi_\omega(t)=(\kappa(t)+\omega)\bmod q$. Then
\[
  \frac1q\sum_{\omega=0}^{q-1}\cost(\Psi_\omega)\;\le\;\frac{\#\{\text{long hops}\}}q+\frac{6T}{(m+1)2^{s_m}}.
\]
In particular, for every $\eps>0$ the walk DAG has a cut of at most $\eps T$ edges and depth at most
$\lceil2/\eps\rceil-1$ as soon as $(m+1)2^{s_m}\ge12/\eps$.
\end{lemma}

\begin{informalproof}
\emph{Hops.} For a hop $(p,u)$ we have $\kappa(u)=\kappa(p)+1$, so $\Psi_\omega(u)\le\Psi_\omega(p)$ if and only if
$\Psi_\omega(p)=q-1$, which holds for exactly one~$\omega$. Averaged over~$\omega$, the number of violated hops is
$\#\{\text{long hops}\}/q$ (all hops of~$U$ are long).

\emph{Descents.} A descent at~$t$ requires $\kappa(t)\ne\kappa(t-1)$; we count such~$t$. Let the step from $t-1$ to~$t$
be emitted by a segment of a level-$i$ stroke~$S$, $i<m$, running from $a_j$ to the grid point~$b_j$ in
$2^{s_{i+1}+1}$ steps, so that its middle cell is $a_j+\sigma2^{s_{i+1}}$. If neither $x(t-1)$ nor $x(t)$ is one of
$a_j$, $a_j+\sigma2^{s_{i+1}}$ and~$b_j$, both cells lie in one open half-segment of~$S$, both times belong to the
segment, and $\kappa(t-1)=\kappa(t)$: indeed, for a segment time~$s$ at such a cell~$c$, the visits to~$c$ before~$s$
are those before the first time of~$S$, because by \cref{lem:dep-structure}(a) the earlier parts of~$S$ stay between
the first cell of~$S$ and~$a_j$ and the segment reaches~$c$ only at~$s$; and the number of earlier visits is the same
for all cells of the open half-segment by \cref{lem:dep-structure}(d). So at most four steps emitted by the segment
change~$\kappa$: those leaving~$a_j$, entering and leaving the middle cell, and entering~$b_j$. In the same way a
level-$m$ stroke from~$a$ to $a+\sigma2^{s_m}$ emits at most two such steps, those leaving~$a$ and entering
$a+\sigma2^{s_m}$. A period contains $4\cdot2^{gi}$ strokes of level~$i$, so the number of changes of~$\kappa$ per
period is at most
\[
  4\sum_{i<m}4\cdot2^{gi}\cdot2^{g-1}+2\cdot4\cdot2^{gm}\;\le\;\frac{8\cdot2^g}{2^g-1}\,2^{gm}+8\cdot2^{gm}
  \;\le\;24\cdot2^{gm}=\frac{6P}{(m+1)2^{s_m}},
\]
using $2^{gm}=2^{s_0-s_m}$. Each change of~$\kappa$ is at most one descent of~$\Psi_\omega$, for every~$\omega$, which
gives the second term.

For the last claim take $q=\lceil2/\eps\rceil$: the first term is at most $\eps T/2$, the second at most $\eps T/2$,
and some~$\omega$ does at least as well as the average; \cref{lem:dep-label} turns $\Psi_\omega$ into a cut.
\end{informalproof}

\paragraph{Why it matters.} One deeply nested walk costs essentially nothing beyond the floor $1/q$ per hop, at every
depth of nesting, and larger~$m$ even helps. The label is a single anchored digit, since $\kappa$ rises by exactly~$1$
along every hop whatever happens in between, and it is not a band rule, so it escapes \cref{rem:dep-N,rem:dep-Nprime}. The
difficulty of the instance studied next lies entirely in the interaction of two heads.

\subsection{The deep-nesting pair}\label{sec:dep-dnp}

\Cref{thm:B} covers walks of linear leg area, in particular walks whose legs all begin or end near a home. In
$U(s_0,m,g)$ most legs do neither: the nesting happens far from home, at every level. Two copies of~$U$ shifted
against each other in time give the instance that we regard as the sharpest open case of \cref{prob:tcp} for $k=2$.

\begin{definition}[The deep-nesting pair]\label{def:dnp}
Let $s_0$, $m$, $g$ be as in \cref{def:dep-walk} and $P=4(m+1)2^{s_0}$. The \emph{golden shift} $\xi$ is the least odd
integer at least $\lfloor0.381966\,P\rfloor$. The \emph{deep-nesting pair} $\mathrm{DNP}(s_0,m,g)$ consists of the
walks $x_1=U(s_0,m,g)$ and $x_2(t)=x_1(t+\xi)$, $t\in\Z$. For $r\ge1$, its \emph{$r$-fold cover} is the pair on the
cyclic time axis $\Z_T$, $T=rP$. For $t\in\Z_T$ and a head~$h$, $\prev_h(t)$ is the residue modulo~$T$ of the previous
visit, on~$\Z$, of head~$h$ to $x_h(t)$; this defines bijections $\prev_1,\prev_2$ of $\Z_T$, and the \emph{hops} of
head~$h$ on the cover are the pairs $(\prev_h(t),t)$, $t\in\Z_T$, so there are $n_{\mathrm h}=2T$ hops in all. A
\emph{label} on the cover is a map $G\colon\Z_T\to\{0,\dots,q-1\}$ with $q\ge1$. A hop $(p,u)$ is \emph{violated} if
$G(u)\le G(p)$, a time~$t$ is a \emph{descent} if $G(t)<G(t-1)$, and $\cost(G)$ is the number of violated hops of both
heads plus the number $\mathrm{desc}(G)$ of descents. Hops are indexed by head and later endpoint, so a pair that is a
hop of both heads would count twice: $\cost(G)$ counts violations per head, unlike $\cost(\Psi)$ on the walk DAG;
by \cref{rem:dep-linear} no such pair exists when $s_m\ge2$.
\end{definition}

The constant $0.381966$ is $2-\phi$ to six decimals, $\phi$ the golden ratio, and we also call $\mathrm{DNP}(s_0,m,g)$
the \emph{golden pair}. It is meant to keep the two hierarchies
misaligned, but because of the truncation it fails to do so when $15625$ divides $m+1$: then $\xi\equiv1\pmod{2^{s_0-3}}$,
and the hierarchies are aligned at every scale up to $2^{s_0-3}$; for some other~$m$ they nearly align at several
scales (\cref{rem:dep-Em}). Only the arithmetic of \cref{thm:E1}, where $m=1$,
depends on its exact value.

The cover is a model of the walk
DAG of many periods in which the identities below are exact; \cref{rem:dep-linear} transfers the main bound back to
the walk DAG. A label on the cover plays the role of a label of depth $q-1$; by \cref{lem:dep-label}, labels lose
nothing on the walk DAG. A cell~$c$ is \emph{interior} if $\abs c<2^{s_0}$; by \cref{lem:dep-structure}(c), on the
$r$-fold cover each head visits every interior cell $r(2m+2)$ times and each of the two tips $\pm2^{s_0}$ $r(m+1)$
times.

For $t\in\Z_T$, identified with its representative in $[0,T)$, the \emph{visit index} $\kappa_h(t)$ is the number of
times $s\in[0,t)$ with $x_h(s)=x_h(t)$. The \emph{error} of a hop $(p,u)$ is
$e(p,u)=(G(u)-G(p)-1)\bmod q\in\{0,\dots,q-1\}$, and $G$ is \emph{exact} on the hop if the error is~$0$, that is,
$G(u)\equiv G(p)+1\pmod q$. The labels $(\kappa_1+\omega)\bmod q$ of \cref{lem:dep-visit} are exact on every head-$1$
hop at an interior cell when $q$ divides $r(2m+2)$. The \emph{level} of a hop $(p,u)$ of head~$h$ is the least level
of a turn of~$x_h$ strictly inside $(p,u)$ (on the cover, inside any lift of the hop to~$\Z$); every hop contains a
turn, its owner, so levels are defined.

\paragraph{Anchored counts.} The visit index is anchored: it rises by exactly one along every hop of its head. A
count restricted to a band of levels rises by exactly one along the hops of that band and is unchanged along the
others.

\begin{lemma}[Level-band anchored counts]\label{lem:dep-band}
\notformalized
For a head~$h$, a set~$B$ of levels and a time $t\ge0$, let $U_h^B(t)$ be the number of times $s\in[0,t)$ with
$x_h(s)=x_h(t)$ such that the hop of head~$h$ that begins at~$s$ has its level in~$B$. Then for every hop $(p,u)$ of
head~$h$ with $p\ge0$, $U_h^B(u)-U_h^B(p)$ is $1$ if the level of $(p,u)$ lies in~$B$ and $0$ otherwise.
\end{lemma}

\begin{informalproof}
The visits to $x_h(u)$ in $[0,u)$ are those in $[0,p)$ and $p$ itself, because $p$ and~$u$ are consecutive visits.
The visit~$p$ is counted in $U_h^B(u)$ exactly when the hop that begins at~$p$, which is $(p,u)$, has its level in~$B$.
\end{informalproof}

With $B$ all levels, $U_h^B$ is the visit index. Along the \emph{chain} of a cell, the visits of one head (here head~$1$) to it
in time order, a label exact on the head-$1$ hops whose
level lies in~$B$ rises by one modulo~$q$ wherever $U_1^B$ rises, so modulo~$q$ it differs from $U_1^B$ by a function
of the cell and of the number of earlier head-$1$ hops at that cell with level outside~$B$; this is the class used in
\cref{rem:dep-Em}.

\begin{remark}[Numerical: how often band counts change]\label{rem:dep-bandcount}
Numerically, on all instances checked, $U_1^{\{0,\dots,j\}}$ is constant during every V
of level greater than~$j$ at all times at which the head is not at the base of the V; at the base it may change, so the
plain statement ``constant during the V'' is false. The total number of its changes was of the order of
$T/((m+1)2^{s_{j+1}})+T(m-j)/((m+1)2^{s_j})$, so the gain $1/(m+1)$ of \cref{lem:dep-visit} is lost for $j<m$.
Nothing below uses this remark.
\end{remark}

\begin{lemma}[No monotone anchoring]\label{lem:dep-monotone}
\notformalized
Let $x$ be a time shift of $U(s_0,m,g)$, let $1\le i<m$, and let $\mathcal V$ be a V of level~$i$ with valley
cell~$v$ whose forward stroke has direction~$\sigma$. For $1\le\rho<2^{s_i}$ let $p_\rho$ be the last visit to
$v+\sigma\rho$ during the back stroke and $u_\rho$ its first visit during the forward stroke; $(p_\rho,u_\rho)$ is a
hop. Let $K$ be an integer function on the span of~$\mathcal V$ with $K(t)\ge K(t-1)$ whenever $t-1$ and~$t$ belong to
the span, and with $K(u')\ge K(p')+1$ for every hop $(p',u')$ of level $i+1$ with both endpoints in the span. Then
\[
  K(u_\rho)-K(p_\rho)\;\ge\;\rho/2^{s_{i+1}}-3/2\qquad(1\le\rho<2^{s_i});
\]
in particular $K(u_\rho)-K(p_\rho)\ge2^g-2$ for $\rho=2^{s_i}-1$.
\end{lemma}

\begin{informalproof}
After its last visit to $v+\sigma\rho$ the back stroke ends at~$v$, so by continuity it stays on the side of~$v$; the
forward stroke starts at~$v$ and first reaches $v+\sigma\rho$ at~$u_\rho$. So $p_\rho,u_\rho$ are consecutive visits.
Put $s=s_{i+1}$ and let $\gamma_k=v+\sigma k2^{s+1}$. By \cref{lem:dep-structure}(a), the back stroke has V's of
level $i+1$ attached at $\gamma_k$ with ranges $[\gamma_k,\gamma_k+\sigma2^s]$, $0\le k<2^{g-1}$, and the forward
stroke has V's attached at $\gamma_{k+1}$ with ranges $[\gamma_{k+1}-\sigma2^s,\gamma_{k+1}]$. If
$(2k+1)2^s<\rho$, the back stroke's V at~$\gamma_k$ lies strictly on the side of~$v$ of $v+\sigma\rho$; it is
attached after the back stroke has passed $v+\sigma\rho$, and nothing after it in the back stroke returns to
$v+\sigma\rho$, so it lies inside $(p_\rho,u_\rho)$. There are at least $(\rho/2^s-1)/2$ such~$k$. If
$(k+1)2^{s+1}<\rho$, the forward stroke reaches $\gamma_{k+1}$ and completes its V there before it reaches
$v+\sigma\rho$, so this V lies inside $(p_\rho,u_\rho)$ as well; there are at least $\rho/2^{s+1}-1$ such~$k$. In all,
at least $\rho/2^s-3/2$ V's of level $i+1$ lie inside $(p_\rho,u_\rho)$, and their spans are disjoint.

In each of them, with valley cell $v'$ and forward direction $\sigma'$, take the hop between the last visit of the back
stroke and the first visit of the forward stroke to $v'+\sigma'$. As $2^s\ge2$, this cell is not the base, so the hop
lies after the apex time; by \cref{lem:dep-structure}(b) the only turns strictly inside it are the valley time of this
V and turns of V's nested in it, so the hop has level $i+1$. Listing the chosen hops in time order and joining them by
the stretches of time between them, $K$ gains at least~$1$ along each chosen hop and does not decrease along the
stretches. Hence $K(u_\rho)-K(p_\rho)$ is at least the number of chosen hops. For $\rho=2^{s_i}-1$ the bound is
$2^g-2^{-s}-3/2\ge2^g-2$.
\end{informalproof}

\paragraph{Clock labels.} Counting turns instead of visits gives digits that never decrease. One digit per scale
suffices for both heads, because the two walks of a deep-nesting pair have the same scales.

\begin{proposition}[Clock labels]\label{prop:dep-clock}
\notformalized
Let $x_1$, $x_2$ be time shifts of $U(s_0,m,g)$, for instance a deep-nesting pair, on $[0,T)$, where $T$ is a multiple
of~$P$. For $0\le i\le m$ let $N^{(i)}(t)$ be the number of turns of level~$i$ of either head at times in $[0,t]$.
For an
integer $B\ge2$ and $\omega\in\{0,\dots,B-1\}^{m+1}$ let
\[
  \Psi_\omega(t)=\sum_{i=0}^{m}B^{m-i}\bigl((N^{(i)}(t)+\omega_i)\bmod B\bigr)\in\{0,\dots,B^{m+1}-1\}.
\]
Then the average of $\cost(\Psi_\omega)$ over all~$\omega$ is at most $12T/B$. Consequently, for every $\eps>0$ the walk
DAG of the pair has a cut of at most $\eps T$ edges and depth at most $\lceil12/\eps\rceil^{m+1}-1$.
\end{proposition}

\begin{informalproof}
The vector $(N^{(0)}(t),\dots,N^{(m)}(t))$ is coordinatewise non-decreasing in~$t$, and $\Psi_\omega(t)$ is the integer
whose base-$B$ digits, most significant first, are $(N^{(i)}(t)+\omega_i)\bmod B$.

\emph{Hops.} Let $(p,u)$ be a hop of head $h'\in\{1,2\}$ inside $[0,T)$. Its owner is a turn strictly inside $(p,u)$,
so some $N^{(i)}$ increases from $p$ to~$u$; let $j$ be the least such~$i$ and $\Delta=N^{(j)}(u)-N^{(j)}(p)\ge1$. The
digits of index less than~$j$ agree at $p$ and~$u$. If $(N^{(j)}(p)+\omega_j)\bmod B<B-\Delta$, digit~$j$ increases and
$\Psi_\omega(u)>\Psi_\omega(p)$ whatever the later digits are; for uniform $\omega_j$ the residue
$(N^{(j)}(p)+\omega_j)\bmod B$ is uniform, so the hop is violated with probability at most $\Delta/B$. Now $\Delta$ is
the number of turns of level~$j$ of either head at times in $(p,u]$, and no turn of level less than~$j$ of either head
lies in $(p,u]$. Fix a turn~$\tau$ of level~$j$ and consider the hops $(p,u)$ of head~$h'$ with $\tau\in(p,u]$ and
without turns of head~$h'$ of level less than~$j$ strictly inside. By \cref{lem:dep-structure}(e) if $j\ge1$, and
because $\abs{x_{h'}}\le2^{s_0}$ if $j=0$, the cell of every such hop lies within distance $2^{s_j}$ (within
$2^{s_0+1}$ if $j=0$) of $x_{h'}(\tau)$; and at each cell at most one hop of head~$h'$ has $\tau\in(p,u]$, because the
hops of one head at one cell are disjoint half-open intervals of time. So the turn~$\tau$ is counted, over both
heads, at most $2(2^{s_j+1}+1)$ times if $j\ge1$ and at most $2(2^{s_0+1}+1)$ times if $j=0$. By
\cref{lem:dep-structure}(b), each head has, per period, two turns of level~$0$ and $4\cdot2^{gj}$ turns of level~$j$
for $1\le j\le m$ (a period contains $2\cdot2^{gj}$ V's of level~$j$). Summing over the turns of both heads in one
period gives at most
\[
  4\cdot2(2^{s_0+1}+1)+\sum_{j=1}^{m}8\cdot2^{gj}\cdot2(2^{s_j+1}+1)
  =32(m+1)2^{s_0}-16\cdot2^{s_0}+8+16\sum_{j=1}^m2^{gj}\;\le\;8P+8\;\le\;9P,
\]
because $16\sum_{j\le m}2^{gj}\le32\cdot2^{gm}\le16\cdot2^{s_0}$ (as $s_m\ge1$) and $P\ge8$. So the expected number of
violated hops is at most $9T/B$.

\emph{Descents.} If no turn occurs at~$t$, then $\Psi_\omega(t)=\Psi_\omega(t-1)$. Otherwise the coordinates increase
by one for each head that turns at~$t$, so by at most two in all; the first changed digit increases by $1$ or~$2$, and
$\Psi_\omega$ descends only if this digit wraps, which has probability at most (number of turns at~$t$)$/B$. The two
heads have at most
$4+8\sum_{j\le m}2^{gj}\le4+8\cdot2^{s_0}\le3P$ turns per period, so the expected number of descents is at most
$3T/B$.

For the last claim we may assume $\eps<12$, since cutting all of the at most $3T$ edges gives depth~$0$; take
$B=\lceil12/\eps\rceil\ge2$, choose~$\omega$ with cost at most the average, and apply
\cref{lem:dep-label}.
\end{informalproof}

\paragraph{Why they matter.} By \cref{prop:dep-clock} every deep-nesting pair of fixed nesting depth~$m$ is depth-reducible,
with $D(\eps)$ of the form $(c/\eps)^{m+1}$; so depth-robustness can only appear as $m$ grows. The two lemmas show why
the obvious ways of beating the clock fail. For $g\ge2$, a digit that decides one head's hops at several levels exactly,
with slack below $2^g-2$, must descend inside every V of that head (\cref{lem:dep-monotone}; for $g=1$ the lemma is
empty), and the other head's long
hops that span such a descent see it. A digit that decides the other head's hops must instead be non-decreasing across
them. Anchored digits (\cref{lem:dep-visit,lem:dep-band}) and clock digits are the two ends of this spectrum, and
every explicit label we know is built from them.

\paragraph{Packing certificates.} A proof that the pairs are depth-robust as $m$ grows must show that every cut of depth~$D$ is
large. The standard way to certify such lower bounds, linear programming duality or packings of local obstructions,
cannot do it.

\begin{proposition}[No packing certificate]\label{prop:packing}
\notformalized
Let $x_1,\dots,x_k$ be walks of length~$T$, let $H$ be the set of long hops of their walk DAG, and let $D\ge0$. Call a
set~$F$ of edges \emph{$D$-forcing} if no map $\Psi\colon[0,T)\to\{0,\dots,D\}$ satisfies $\Psi(a)\ge\Psi(b)+1$ for
every long hop $a\to b$ in~$F$ and $\Psi(a)\ge\Psi(b)$ for every state edge $a\to b$ in~$F$.
\begin{enumerate}[label=(\alph*)]
\item Every $D$-forcing set contains a path with at least $D+1$ long hops.
\item Let $F_1,\dots,F_n$ be $D$-forcing and $y_1,\dots,y_n\ge0$ with $\sum_{j:e\in F_j}y_j\le1$ for every edge~$e$.
Then every cut of depth at most~$D$ has at least $\sum_jy_j$ edges, and $\sum_jy_j\le\card H/(D+1)\le kT/(D+1)$.
\item The linear program that minimises $\sum_ez_e$ over $z\ge0$ subject to $\sum_{e\in\pi}z_e\ge1$ for every path~$\pi$
with at least $D+1$ long hops has value at most $\card H/(D+1)$.
\end{enumerate}
Hence a lower bound $\eps T$ on the size of all cuts of depth at most~$D$ can be certified in the form~(b) or~(c) only
if $D<k/\eps$; such certificates never show depth-robustness.
\end{proposition}

\begin{informalproof}
(a) If $F$ contains no path with $D+1$ long hops, let $\Psi(a)$ be the largest number of long hops on a path in~$F$
that starts at~$a$. Then $0\le\Psi\le D$, and for every edge $a\to b$ of~$F$, $\Psi(a)\ge\Psi(b)+1$ if it is a long
hop and $\Psi(a)\ge\Psi(b)$ otherwise; so $F$ is not forcing.

(b) Let $F$ be a cut of depth at most~$D$, and let $\Psi$ be the map of the proof of \cref{lem:dep-label}; it satisfies
the constraints of all edges outside~$F$. Since each $F_j$ is forcing, it contains an edge of~$F$, and
$\card F\ge\sum_{e\in F}\sum_{j:e\in F_j}y_j\ge\sum_jy_j$. By~(a), $\card{F_j\cap H}\ge D+1$, so
$(D+1)\sum_jy_j\le\sum_jy_j\card{F_j\cap H}=\sum_{e\in H}\sum_{j:e\in F_j}y_j\le\card H$. Finally $\card H\le kT$,
because each time is the later endpoint of at most one hop of each head.

(c) Put $z_e=1/(D+1)$ on long hops and $z_e=0$ on state edges.

For the last sentence, $\card H/(D+1)\ge\eps T$ forces $D+1\le k/\eps$.
\end{informalproof}

On a cyclic time axis, such as the covers of \cref{def:dnp}, every cyclic hop chain is forcing: a chain of $n$ hops
always contains at least $\lceil n/q\rceil$ violated hops, more than the average $n/q$ unless $q$ divides~$n$. This is
why we use covers on which $q$ divides the visit numbers of the cells that matter; the effect does not arise on the
walk DAG, which is acyclic.

\paragraph{Why it matters.} \Cref{prop:packing} rules out the natural strategy for proving depth-robustness by duality:
packings of forcing structures, of any shape, certify at most what the trivial assignment of weight $1/(D+1)$ to
every long hop certifies. \Cref{lem:dep-euler}(c) below records the same barrier for cycle bounds in the cost
identity. A robustness proof cannot be a fractional packing; it has to exploit integrality, and we expect it to be
multiplicative: it would show that a label with $q$ values cannot serve many levels of both heads at once.

\subsection{Exact identities for labels}\label{sec:dep-identities}

On the cover every label obeys an exact cost identity, which separates the unavoidable floor $1/q$ per hop from the
errors of the hops.

\begin{lemma}[Lemma Zq: the cost identity]\label{lem:Zq}
\formal{DepthLabels.lean}{lemma_Zq, card_violated_eq, card_descents_eq}
Let $\tau$ and $H$ be finite sets, $q\ge1$ an integer, $G\colon\tau\to\{0,\dots,q-1\}$, and $t\mapsto t^-$ a bijection
of~$\tau$. Let $p,u\colon H\to\tau$ be maps and $\mathrm{succ}$ a bijection of~$H$ with
$p(\mathrm{succ}(h))=u(h)$ for all $h\in H$. Put $e_h=(G(u(h))-G(p(h))-1)\bmod q$ and $\chi_t=(G(t)-G(t^-))\bmod q$,
residues in $\{0,\dots,q-1\}$. Then
\[
  q\cdot\#\{h:G(u(h))\le G(p(h))\}=\card H+\sum_he_h,\qquad q\cdot\#\{t:G(t)<G(t^-)\}=\sum_t\chi_t,
\]
and hence $q\,(\#\{h:G(u(h))\le G(p(h))\}+\#\{t:G(t)<G(t^-)\})=\card H+\sum_he_h+\sum_t\chi_t$.
\end{lemma}

\begin{outline}
For $-q\le y<q$ one has $y\bmod q=y+q\cdot[y<0]$ (\lean{emod_of_bounds}). With $y=G(u(h))-G(p(h))-1$, which lies in
$[-q,q-2]$, this gives $e_h=G(u(h))-G(p(h))-1+q\cdot[G(u(h))\le G(p(h))]$. Summing over~$H$ and using
$\sum_hG(u(h))=\sum_hG(p(h))$, obtained by reindexing along the bijection $\mathrm{succ}$, gives the first identity. The
second is the same computation with $y=G(t)-G(t^-)\in(-q,q)$ and $\sum_tG(t^-)=\sum_tG(t)$; the third is their sum.
\end{outline}

On the cover of \cref{def:dnp} take $\tau=\Z_T$, $t^-=t-1$, let $H=\{1,2\}\times\Z_T$ index the hops of the two heads, with
$p,u$ their endpoints, and let $\mathrm{succ}$ map a hop to the next hop of the same head at the same cell; the hypothesis only
says that the hops of each head at each cell form a cycle. Since $\sum_t\chi_t=q\,\mathrm{desc}(G)$, the lemma reads
\begin{equation}\label{eq:dep-Zq}
  \cost(G)\;=\;\mathrm{desc}(G)+\frac1q\sum_{\text{hops }h}(1+e_h).
\end{equation}
Every hop costs at least $1/q$, and beyond that exactly its error divided by~$q$; every descent costs~$1$. For $q=2$,
\eqref{eq:dep-Zq} says that twice the cost is $n_{\mathrm h}$ plus the number of hops with $G(u)=G(p)$ plus the number
of times with $G(t)\ne G(t-1)$, so depth~$1$ is a problem of frustration in a signed graph. The identity also shows
that $\cost(G)$ does not change when $G$ is replaced by $(G+c)\bmod q$.

\begin{corollary}[Per-class split]\label{cor:dep-split}
\notformalized
On the cover of a deep-nesting pair, for every label~$G$,
\[
  \cost(G)=\mathrm{desc}(G)+\sum_{h=1}^{2}\sum_{i=0}^{m}\frac1q\sum_{\substack{\text{hops of head }h\\\text{of level
}i}}
  (1+e).
\]
\end{corollary}

\begin{informalproof}
Group the hops in \eqref{eq:dep-Zq} by head and level; each hop has exactly one of each.
\end{informalproof}

So the classes (head, level) contribute additively, and how well $G$ serves a class is measured by the error mass of
the class: a class on which $G$ is exact costs its size over~$q$.

\paragraph{Lifts.} To state the next lemmas without boundary terms we lift labels, walks and hops to~$\Z$: $G$ and the
walks are $T$-periodic, the hops of head~$h$ are the pairs $(\prev_h(u),u)$, $u\in\Z$, with $\prev_h$ the previous
visit on~$\Z$, and each hop of the cover is the residue of such a pair. For a head~$h$, a cell~$c$ and a visit
$s\in\Z$ of head~$h$ to~$c$, let $\hat\kappa_h(s)$ be the number of visits of head~$h$ to~$c$ at times in $[0,s)$ if
$s\ge0$, and minus the number at times in $[s,0)$ if $s<0$. Thus $\hat\kappa_h(s')-\hat\kappa_h(s)=1$ for consecutive
visits $s<s'$, $\hat\kappa_h(s+T)=\hat\kappa_h(s)+n_h(c)$, where $n_h(c)$ is the number of visits of head~$h$ to~$c$ per
cover, and $\hat\kappa_h=\kappa_h$ on $[0,T)$. Let $t_0(c)$ be the first visit of head~$h$ to~$c$ in $[0,T)$, and let
$\hat E_h(s)$ be the total error of the hops of head~$h$ at~$c$ whose later endpoint lies in $(t_0(c),s]$ if
$s\ge t_0(c)$, and minus the total error of those with later endpoint in $(s,t_0(c)]$ if $s<t_0(c)$. So
$\hat E_h(s')-\hat E_h(s)\ge0$ is the error mass of the hops of~$c$ ending in $(s,s']$ when $s<s'$.

\begin{lemma}[Chain decomposition]\label{lem:chain}
\partformal{DepthLabels.lean}{congr_of_exact}{the general form~(a), with error terms, and its reading~(c) on the cover.}
\begin{enumerate}[label=(\alph*)]
\item For every head~$h$, every cell~$c$ and every visit $s\in\Z$ of head~$h$ to~$c$,
\[
  G(s)\equiv G(t_0(c))+\hat\kappa_h(s)+\hat E_h(s)\pmod q,
\]
and the total error of the $n_h(c)$ hops of head~$h$ at~$c$ on the cover is congruent to $-n_h(c)$ modulo~$q$.
\item (Exact case, abstract form.) Let $\prev\colon\tau\to\tau$ and $\mathrm{cell}\colon\tau\to C$ be maps on a set
$\tau$ with $\mathrm{cell}(\prev(t))=\mathrm{cell}(t)$, such that whenever $\mathrm{cell}(s)=\mathrm{cell}(t)$ some
iterate of $\prev$ maps $t$ to~$s$. Let $S\subseteq C$, and let $G,\kappa\colon\tau\to\Z$ satisfy
$G(t)-G(\prev(t))\equiv1$ and $\kappa(t)-\kappa(\prev(t))\equiv1\pmod q$ whenever $\mathrm{cell}(t)\in S$. Then
$G(s)-\kappa(s)\equiv G(t)-\kappa(t)\pmod q$ whenever $\mathrm{cell}(s)=\mathrm{cell}(t)\in S$.
\item If $G$ is exact on all hops of head~$h$ at a cell~$c$, then $q$ divides $n_h(c)$ and
$G\equiv\kappa_h+G(t_0(c))\pmod q$ at the visits of~$c$ in $[0,T)$. Conversely, if $q$ divides $n_h(c)$ and
$G-\kappa_h$ is constant modulo~$q$ on the visits of~$c$ in $[0,T)$, then $G$ is exact on all hops of head~$h$ at~$c$.
\end{enumerate}
\end{lemma}

\begin{outline}[Outline of the formal proof of~(b)]
Induction on~$n$ shows $G(\prev^n(t))-\kappa(\prev^n(t))\equiv G(t)-\kappa(t)$: the cell stays in~$S$ because
$\mathrm{cell}\circ\prev=\mathrm{cell}$, and at each step the two congruences cancel. The orbit hypothesis supplies~$n$
with $\prev^n(t)=s$.
\end{outline}

\begin{proof}[Proof of~(a) and~(c) (not formalized)]
(a) For consecutive visits $s<s'$ of head~$h$ to~$c$, $G(s')-G(s)\equiv1+e(s,s')$ by the definition of the error,
while $\hat\kappa_h$ rises by~$1$ and $\hat E_h$ by $e(s,s')$. So $G-\hat\kappa_h-\hat E_h$ is constant modulo~$q$ on
the visits of~$c$, and at $t_0(c)$ it equals $G(t_0(c))$. At $s=t_0(c)+T$ we have $G(s)=G(t_0(c))$,
$\hat\kappa_h(s)=n_h(c)$, and $\hat E_h(s)$ is the total error of the $n_h(c)$ hops of~$c$ on the cover; this gives the
second claim.

(c) If all these errors vanish, (a) gives $0\equiv-n_h(c)$ and $G(s)\equiv G(t_0(c))+\kappa_h(s)$ for
$s\in[0,T)$. Conversely, let $t_0<\dots<t_{n-1}$ be the visits of~$c$ in $[0,T)$. For $1\le j<n$ the hop
$(t_{j-1},t_j)$ satisfies $G(t_j)-G(t_{j-1})\equiv\kappa_h(t_j)-\kappa_h(t_{j-1})=1$, and the remaining hop of~$c$ on
the cover, from $t_{n-1}$ to~$t_0$, satisfies $G(t_0)-G(t_{n-1})\equiv0-(n-1)\equiv1$ because $q$ divides~$n$.
\end{proof}

The visit numbers of the tips are $r(m+1)$ and those of interior cells $r(2m+2)$. So if $q$ divides $r(2m+2)$ but not
$r(m+1)$, for instance $r=q/\gcd(q,4)$ with $4\mid q$ at $m=1$, no label is exact on all head-$1$ hops; the natural
class is that of labels exact on all head-$1$ hops \emph{at interior cells}, which by~(c) are the labels
$G\equiv\kappa_1+\Phi(x_1)\pmod q$ at interior cells, for arbitrary functions~$\Phi$ of the cell. It contains the
visit-index labels of \cref{lem:dep-visit}.

\begin{lemma}[Transfer identity]\label{lem:transfer}
\partformal{DepthLabels.lean}{cycleDefect, cycle_errors_congr}{the general form~(a), with head-$1$ error mass, and the
absorption reading~(c).}
\begin{enumerate}[label=(\alph*)]
\item Let $h_j=(p_j,u_j)$, $j\in\Z_\ell$, be hops of head~$2$ on~$\Z$ with $x_1(u_j)=x_1(p_{j+1})$ for every~$j$. Put
$\Delta\hat\kappa_1(h_j)=\hat\kappa_1(u_j)-\hat\kappa_1(p_j)$ and define the \emph{defect}
$\delta=\sum_j(\Delta\hat\kappa_1(h_j)-1)$. Then
\[
  \sum_je(h_j)\;\equiv\;\delta-\sum_j\bigl(\hat E_1(p_{j+1})-\hat E_1(u_j)\bigr)\pmod q.
\]
\item (Exact case, abstract form.) Let $\prev_2\colon\tau\to\tau$, $\mathrm{cell}\colon\tau\to C$, $S\subseteq C$ and
$G,\kappa\colon\tau\to\Z$ be such that $G-\kappa$ is constant modulo~$q$ on the times of each cell of~$S$. Let
$b\colon\Z_\ell\to\tau$ satisfy $\mathrm{cell}(b(j))=\mathrm{cell}(\prev_2(b(j+1)))$ and $\mathrm{cell}(b(j))\in S$ for
all~$j$. Then
\[
  \sum_j\bigl(G(b(j))-G(\prev_2(b(j)))-1\bigr)\;\equiv\;\sum_j\bigl(\kappa(b(j))-\kappa(\prev_2(b(j)))-1\bigr)
  \pmod q;
\]
the right side is the defect \lean{cycleDefect} of $\kappa$, $\prev_2$ and~$b$.
\item (Absorption.) In~(a), call the head-$1$ stretch from $u_j$ to $p_{j+1}$ \emph{forward} if $u_j<p_{j+1}$ and
\emph{backward} otherwise, and let $A_j=\abs{\hat E_1(p_{j+1})-\hat E_1(u_j)}$ be its error mass. If $-q<\delta<0$ and
$\sum_je(h_j)<q-\abs\delta$, then either the backward stretches carry error mass at least
$\abs\delta+\sum_je(h_j)$, or the forward stretches carry error mass at least
$q-\abs\delta-\sum_je(h_j)+\sum_{\text{backward}}A_j$.
\end{enumerate}
\end{lemma}

\begin{outline}[Outline of the formal proof of~(b)]
After reindexing one of the sums by the rotation $j\mapsto j+1$ of $\Z_\ell$, the difference of the
two sides becomes $\sum_j[(G-\kappa)(b(j))-(G-\kappa)(\prev_2(b(j+1)))]$. Each term is divisible by~$q$, because
$b(j)$ and $\prev_2(b(j+1))$ are times of the same cell of~$S$.
\end{outline}

\begin{proof}[Proof of~(a) and~(c) (not formalized)]
(a) By the definition of the error, $G(u_j)-G(p_j)\equiv1+e(h_j)$. By \cref{lem:chain}(a) for head~$1$ at the cell
$x_1(u_j)=x_1(p_{j+1})$, whose base term cancels,
$G(p_{j+1})-G(u_j)\equiv\hat\kappa_1(p_{j+1})-\hat\kappa_1(u_j)+\hat E_1(p_{j+1})-\hat E_1(u_j)$. Summing both
congruences over $j\in\Z_\ell$, the values of~$G$ cancel, and $\sum_j(\hat\kappa_1(p_{j+1})-\hat\kappa_1(u_j))
=-\sum_j\Delta\hat\kappa_1(h_j)$. Rearranging gives~(a).

(c) Write $\sigma_e=\sum_je(h_j)$, $\delta'=-\delta\in(0,q)$, and let $F_A$ and $B_A$ be the error masses of the forward
and of the backward stretches. The term $\hat E_1(p_{j+1})-\hat E_1(u_j)$ is $A_j$ for a forward and $-A_j$ for a
backward stretch, so~(a) says $B_A-F_A\equiv\sigma_e+\delta'\pmod q$, where $0<\sigma_e+\delta'<q$. Write
$B_A-F_A=\sigma_e+\delta'+qz$ with $z\in\Z$. If $z\ge0$, then $B_A\ge\sigma_e+\delta'$. If $z\le-1$, then
$F_A=B_A-\sigma_e-\delta'-qz\ge q-\delta'-\sigma_e+B_A$.
\end{proof}

For labels exact on head~$1$ at interior cells, all errors of head-$1$ hops at interior cells vanish, so by~(a) or~(b)
every closed sequence of head-$2$ hops through interior head-$1$ cells has $\sum_je(h_j)\equiv\delta$: a cell
correction~$\Phi$ telescopes and absorbs nothing. In general, by~(c), a defect $-q<\delta<0$ can be absorbed cheaply,
at a cost of order $\abs\delta/q$, only by errors of head-$1$ hops on backward stretches, which lie inside the overlaps of
consecutive head-$2$ hops; errors on forward stretches enter with the wrong sign and help only after a wrap, at a cost of nearly a whole
unit.

\begin{remark}[Numerical: who absorbs in optimised labels]\label{rem:dep-absorb}
For the deep-nesting pair with $m=3$, $g=1$ and $s_0=6$, we examined the best labels found by optimisation at
$D=7,15,31,63$; they cost about four to five times less than clock labels with the same number of values. For the
pairs of head-$2$ hops $a\to b$, $b\to a$ in head-$1$ cell space whose defect a label exact on head~$1$ could not
absorb, head-$1$ errors on backward stretches did the absorbing in $55\%$, $83\%$, $88\%$ and $96\%$ of the cases,
and head~$2$ paid (error sum at least $q/2$) in $28\%$, $9\%$, $9\%$ and $4\%$. The transfer identity held on all
$6925$ pairs examined.
\end{remark}

\begin{lemma}[Euler: balanced degrees, chain defects, and the barrier in $\Z_q$ form]\label{lem:dep-euler}
\notformalized
On the $r$-fold cover of a deep-nesting pair:
\begin{enumerate}[label=(\alph*)]
\item The multigraph $M_2$ on the cells of head~$1$ with one edge $x_1(p)\to x_1(u)$ for every hop $(p,u)$ of head~$2$
is balanced: the in-degree and the out-degree of a cell~$c$ both equal the number of visits of head~$1$ to~$c$.
\item Let head~$2$ visit the cell~$c'$ at the times $t_0<\dots<t_{n-1}$ of $[0,T)$. Its hops at~$c'$, lifted to
$(t_{j-1},t_j)$ for $1\le j<n$ and $(t_{n-1},t_0+T)$, satisfy the hypothesis of \cref{lem:transfer}(a), and their
defect is $n_1(x_1(t_0))-n_2(c')$. It is $0$ if $c'$ and $x_1(t_0)$ are both interior cells or both tips.
\item In the setting of \cref{lem:transfer}(a), with $A_j$ as in \cref{lem:transfer}(c),
$\sum_je(h_j)+\sum_jA_j\ge\lVert \delta\rVert_q$, where $\lVert
\cdot\rVert_q$
is the distance to $q\Z$, and $\abs\delta$ is at most the number of hops on the closed walk formed by the~$h_j$ and the
head-$1$ stretches, counted with multiplicity. Consequently, if such closed walks $C_1,\dots,C_n$ use every hop of
either head at most once in all, the bound they yield through \eqref{eq:dep-Zq},
$\cost(G)\ge(n_{\mathrm h}+\sum_i\lVert \delta(C_i)\rVert_q)/q$, is at most $2n_{\mathrm h}/q$.
\end{enumerate}
\end{lemma}

\begin{informalproof}
(a) Since $\prev_2$ is a bijection of $\Z_T$, every time~$t$ is the later endpoint of exactly one hop of head~$2$ and
the earlier endpoint of exactly one; the first enters $x_1(t)$ in~$M_2$ and the second leaves it.

(b) Consecutive hops of the list share an endpoint in time, so the head-$1$ cells match, and the last hop ends at
$t_0+T$, where head~$1$ is at $x_1(t_0)$. The sum of the $\Delta\hat\kappa_1$ telescopes to
$\hat\kappa_1(t_0+T)-\hat\kappa_1(t_0)=n_1(x_1(t_0))$, and there are $n=n_2(c')$ hops. As $x_2$ is a time shift
of~$x_1$, $n_2(c')=n_1(c')$, which is $r(2m+2)$ for interior cells and $r(m+1)$ for tips.

(c) By \cref{lem:transfer}(a), $\delta$ is congruent to $\sum_je(h_j)\pm A_1\pm\dots\pm A_\ell$, whose absolute value
is at most $\sum_je(h_j)+\sum_jA_j$. Next, $\hat\kappa_1(p_{j+1})-\hat\kappa_1(u_j)$ is plus or minus the number of
head-$1$ hops on the $j$-th stretch, so $\delta=-\ell-\sum_j(\hat\kappa_1(p_{j+1})-\hat\kappa_1(u_j))$ and $\abs\delta$
is at most the number of hops on the closed walk. If the walks $C_i$ use every hop at most once, the errors they involve
are distinct terms of $\sum_he_h$, so \eqref{eq:dep-Zq} gives $q\cost(G)\ge n_{\mathrm h}+\sum_he_h\ge
n_{\mathrm h}+\sum_i\lVert \delta(C_i)\rVert_q$; and $\sum_i\lVert \delta(C_i)\rVert_q\le\sum_i\abs{\delta(C_i)}\le
n_{\mathrm h}$.
\end{informalproof}

\paragraph{Why they matter.} \Cref{lem:Zq} makes the depth problem on the cover an exact accounting problem: a label
pays $1/q$ per hop, its errors, and its descents, and the error of a hop is a residue modulo~$q$. \Cref{lem:chain}
says that a label is the visit index of either head plus a cell correction plus that head's integrated errors, and
\cref{lem:transfer} says how the two descriptions interact: around a closed walk of head-$2$ hops, the defect of
head~$1$'s visit index must be paid by head~$2$ or absorbed by head~$1$'s errors on backward stretches. By
\cref{lem:dep-euler}(b), the chain of a head-$2$ cell has zero defect unless exactly one of that cell and the head-$1$
cell of its first visit is a tip, and by \cref{lem:dep-euler}(c) lower bounds obtained from hop-disjoint cycles in
this way are bounded by twice the floor: the identity alone does not give robustness, but it says exactly
what a proof has to count.

\subsection{Labels exact on one head}\label{sec:dep-exact}

The one-head mechanism of \cref{lem:dep-visit} counts exactly along one head; adding an arbitrary correction per cell
gives the natural class of labels exact on that head. For this class the transfer identity has no absorption, and
negative cycles force a price that does not depend on~$q$.

\begin{theorem}[Theorem E: labels exact on one head, abstract form]\label{thm:E}
\formal{DepthLabels.lean}{cost_ge_of_negative_cycles, le_of_modEq_neg}
\begin{enumerate}[label=(\alph*)]
\item If $q\ge1$, $E\ge0$ and $\delta<0$ are integers with $E\equiv\delta\pmod q$, then $E\ge q+\delta$.
\item Let $\tau$ be a finite set, $q\ge1$, $G\colon\tau\to\{0,\dots,q-1\}$, and let $t\mapsto t^-$, $\prev_1$ and
$\prev_2$ be bijections of~$\tau$. Let $\mathrm{cell}\colon\tau\to C$, $S\subseteq C$ and $\kappa\colon\tau\to\Z$ be
such that $G-\kappa$ is constant modulo~$q$ on the times of each cell of~$S$. Let $I$ be a finite set and, for $i\in I$,
let $b_i\colon\Z_{\ell_i}\to\tau$ be maps such that $(i,j)\mapsto b_i(j)$ is injective,
$\mathrm{cell}(b_i(j))=\mathrm{cell}(\prev_2(b_i(j+1)))$ and $\mathrm{cell}(b_i(j))\in S$ for all $i,j$, and
\[
  \delta_i:=\sum_{j\in\Z_{\ell_i}}\bigl(\kappa(b_i(j))-\kappa(\prev_2(b_i(j)))-1\bigr)<0\qquad(i\in I).
\]
Then
\begin{multline*}
  2\card\tau+\sum_{i\in I}(q+\delta_i)\\
  \le\;q\Bigl(\#\{t:G(t)\le G(\prev_1t)\}+\#\{t:G(t)\le G(\prev_2t)\}+\#\{t:G(t)<G(t^-)\}\Bigr).
\end{multline*}
\end{enumerate}
\end{theorem}

\begin{outline}
(a) Write $E-\delta=qz$; then $z\ge1$ because $E\ge0>\delta$, so $E\ge q+\delta$. (b) Apply \lean{card_violated_eq} to
each head, with the hops $(\prev_ht,t)$, $t\in\tau$, and the successor bijection $\prev_h^{-1}$, and apply
\lean{card_descents_eq}; all reduced errors are nonnegative. For each~$i$, the reduced errors of the hops
$(\prev_2(b_i(j)),b_i(j))$ have a sum congruent to the unreduced sum, which is congruent to $\delta_i$ by
\lean{cycle_errors_congr}; so by~(a) it is at least $q+\delta_i$. By injectivity these hops are distinct, so the sum
over~$i$ is at most the total reduced error of head~$2$. Linear arithmetic combines the three identities.
\end{outline}

On the cover of a deep-nesting pair take $\tau=\Z_T$, $t^-=t-1$, the hop maps $\prev_1,\prev_2$, $\mathrm{cell}=x_1$,
$S$ the set of interior cells and $\kappa=\kappa_1$; the right side is then $q\cost(G)$, and $2\card\tau=n_{\mathrm h}$.
If $q$ divides $r(2m+2)$, every label exact on all head-$1$ hops at interior cells satisfies the hypothesis on
$G-\kappa_1$, by \cref{lem:chain}(c). A family $(b_i)$ is a set of hop-disjoint closed walks of head-$2$ hops in~$M_2$
through interior head-$1$ cells, each with negative defect. A closed walk with $-q<\delta_i<0$ forces head-$2$ error at
least $q+\delta_i$, nearly~$q$, on its hops, which by \eqref{eq:dep-Zq} costs nearly one unit, whatever~$q$ is. The
formal statement needs only $\delta_i<0$; for $\delta_i\le-q$ the contribution $q+\delta_i$ is not positive and lowers
the bound, so such walks should be left out of the family. Here $\delta_i$ is computed from the residues~$\kappa$ on the
cover, as in the Lean statement; it differs from the lifted defect of \cref{lem:transfer}(a) by $n_1(\cdot)$ on every
hop that wraps around the cover, and in Theorem~E$_1$ no hop wraps.

\paragraph{Why it matters.} \Cref{prop:packing} and \cref{lem:dep-euler}(c) bound every packing certificate for all
labels by about twice the floor. Theorem~E escapes this because the class is restricted: exactness on one head removes
all absorption (\cref{lem:transfer}), so each negative cycle costs about one unit at every depth. A packing of such
cycles with positive density therefore gives a price per hop independent of~$D$, which rules out the whole class as a
depth-reduction mechanism for the instance; the next theorem supplies such a packing in closed form.

\begin{theorem}[Theorem E$_1$]\label{thm:E1}
\notformalized
Let $s_0\ge4$, $\Lambda=2^{s_0}$, let $x_1=U(s_0,1,1)$, of period $P=8\Lambda$, and let $x_2(t)=x_1(t+\xi)$ with
$\xi=3\Lambda+\eta$ for an integer $\eta$ with $0<8\eta<\Lambda$. Let $r,q\ge1$ with $q\mid4r$. Then every label~$G$ on
the $r$-fold cover (defined as in \cref{def:dnp} for this shift) that is exact on all hops of head~$1$ at interior
cells satisfies
\[
  \cost(G)\;\ge\;\nu_1n_{\mathrm h},\qquad\nu_1=\frac{\Lambda/2-4\eta-1}{16\Lambda},
\]
and $\cost(G)\ge\nu_1n_{\mathrm h}+(1-4\nu_1)n_{\mathrm h}/q$ if $q>4$. For the golden shift of \cref{def:dnp} the
condition $0<8\eta<\Lambda$ holds for every $s_0\ge4$; for $s_0\ge5$, $\nu_1\ge67/4096\approx0.01636$, with equality
exactly at $s_0=8$, and for $s_0=4$, $\nu_1=3/256$.
\end{theorem}

\begin{informalproof}
\emph{The walk.} By \cref{def:dep-walk} with $m=g=1$, each level-$0$ stroke is one segment of $\Lambda$ steps followed
by a V of two plain strokes of $\Lambda/2$ steps. So on $[0,4\Lambda]$
\[
  x_1(t)=\begin{cases}
  t,&0\le t\le\Lambda,\\
  2\Lambda-t,&\Lambda\le t\le3\Lambda/2,\\
  t-\Lambda,&3\Lambda/2\le t\le2\Lambda,\\
  3\Lambda-t,&2\Lambda\le t\le3\Lambda,\\
  t-3\Lambda,&3\Lambda\le t\le7\Lambda/2,\\
  4\Lambda-t,&7\Lambda/2\le t\le4\Lambda,
  \end{cases}
\]
and $x_1(t+4\Lambda)=-x_1(t)$ for all~$t$. In $[0,8\Lambda)$, a cell~$c$ with $0<c<\Lambda/2$ is visited at the times
$c$, $3\Lambda-c$, $3\Lambda+c$, $4\Lambda-c$, and a cell with $\Lambda/2<c<\Lambda$ at $c$, $2\Lambda-c$,
$\Lambda+c$, $3\Lambda-c$; we call $0,1,2,3$ the \emph{indices} of these visits. A cell~$c$ with $-\Lambda<c<0$ is
visited at the times $4\Lambda+t$, for the visits $t$ of $-c$ just listed, with the same indices. Each of these cells
is interior and has $4$ visits per period; so for a time~$t$ of the $k$-th period copy $[8\Lambda k,8\Lambda(k+1))$ of
the cover at such a cell, $\kappa_1(t)$ is $4k$ plus the index of the visit. A pair $(p,u)$ is a hop of head~$2$ if and
only if $(p+\xi,u+\xi)$ is a hop of head~$1$.

\emph{The family.} For every integer~$a$ with $1\le a\le\Lambda/2-4\eta-1$ consider the four head-$2$ hops of the
following table, where the third column is the head-$1$ hop they come from, with its cell, and the last two columns
give the cells of head~$1$ at the two endpoints and the difference of the indices of these visits.
\begin{center}
\small
\begin{tabular}{@{}llllc@{}}
\toprule
& head-$2$ hop $(p,u)$ & head-$1$ hop $(p+\xi,u+\xi)$; its cell & $x_1(p)\to x_1(u)$ & $\Delta\kappa_1$\\
\midrule
$e_1$ & $(a,\Lambda-a-2\eta)$ & $(3\Lambda+a+\eta,4\Lambda-a-\eta)$; $a+\eta$ & $a\to\Lambda-a-2\eta$ & $0$\\
$e_2$ & $(2\Lambda+a+2\eta,3\Lambda-a-4\eta)$ & $(5\Lambda+a+3\eta,6\Lambda-a-3\eta)$; $a+3\eta-\Lambda$
& $\Lambda-a-2\eta\to a+4\eta$ & $-2$\\
$e_3$ & $(3\Lambda-a-4\eta,3\Lambda+a+2\eta)$ & $(6\Lambda-a-3\eta,6\Lambda+a+3\eta)$; $a+3\eta-\Lambda$
& $a+4\eta\to a+2\eta$ & $1$\\
$e_4$ & $(3\Lambda-a-2\eta,3\Lambda+a)$ & $(6\Lambda-a-\eta,6\Lambda+a+\eta)$; $a+\eta-\Lambda$
& $a+2\eta\to a$ & $1$\\
\bottomrule
\end{tabular}
\end{center}
The bounds on~$a$ give $0<a+\eta<a+2\eta<a+3\eta<a+4\eta<\Lambda/2$, which places every cell below in the stated range.
Row~$e_1$: the cell $a+\eta$ lies in $(0,\Lambda/2)$ and is visited at $3\Lambda+(a+\eta)$ and $4\Lambda-(a+\eta)$ with
indices $2$ and~$3$, consecutively; at the times $a$ and $\Lambda-a-2\eta$ head~$1$ is on its first stroke, at cells
$a\in(0,\Lambda/2)$ and $\Lambda-a-2\eta\in(\Lambda/2,\Lambda)$, both with index~$0$. Row~$e_2$: the cell
$\Lambda-a-3\eta$ lies in $(\Lambda/2,\Lambda)$ and is visited at $\Lambda+a+3\eta$ and $2\Lambda-a-3\eta$ with indices
$1$ and~$2$, so its negative $a+3\eta-\Lambda$ is visited at the times of the table with indices $1$ and~$2$,
consecutively; at $2\Lambda+a+2\eta=3\Lambda-(\Lambda-a-2\eta)$ head~$1$ visits
$\Lambda-a-2\eta\in(\Lambda/2,\Lambda)$ with index~$3$, and
at $3\Lambda-(a+4\eta)$ it visits $a+4\eta\in(0,\Lambda/2)$ with index~$1$. Row~$e_3$: the visits of
$a+3\eta-\Lambda$ at $6\Lambda-a-3\eta$ and $6\Lambda+a+3\eta=4\Lambda+3\Lambda-(\Lambda-a-3\eta)$ have indices
$2$ and~$3$; at $3\Lambda+(a+2\eta)$ head~$1$ visits $a+2\eta\in(0,\Lambda/2)$ with index~$2$. Row~$e_4$: the cell
$\Lambda-a-\eta$ lies in $(\Lambda/2,\Lambda)$ and is visited at $2\Lambda-a-\eta$ and $2\Lambda+a+\eta$ with indices
$2$ and~$3$, so its negative $a+\eta-\Lambda$ is visited at the times of the table with indices $2$ and~$3$; at
$3\Lambda-(a+2\eta)$ head~$1$ visits $a+2\eta$ with index~$1$, and at $3\Lambda+a$ it visits $a$ with index~$2$. In
each row $p<u$, and all eight times lie in $[0,4\Lambda)$, so the
differences of~$\kappa_1$ in any period copy are the differences of the indices.

\emph{Closure and defect.} The head-$1$ cells run $a\to\Lambda-a-2\eta\to a+4\eta\to a+2\eta\to a$, so
$x_1(u_j)=x_1(p_{j+1})$ for $j\in\Z_4$, identically in $a$, $\Lambda$ and~$\eta$, and all four cells are interior. The
defect is $(0-1)+(-2-1)+(1-1)+(1-1)=-4$. (At $a=0$ and at $a=\Lambda/2-4\eta$ other negative cycles arise, with other
index patterns; we do not use them.)

\emph{The mirror family.} Since $x_1(t+4\Lambda)=-x_1(t)$ for all~$t$, also $x_2(t+4\Lambda)=-x_2(t)$, so adding
$4\Lambda$ to all times of a member gives four head-$2$ hops whose head-$1$ cells are negated. The cycle condition is
preserved, the indices are the same by the description of the visits of negative cells, and the defect is again~$-4$.

\emph{Disjointness.} A hop of head~$2$ on the cover is determined by its later endpoint. In $[0,8\Lambda)$ the later
endpoints of $e_1$, $e_2$, $e_3$, $e_4$ lie in $(\Lambda/2,\Lambda)$, $(5\Lambda/2,3\Lambda)$, $(3\Lambda,7\Lambda/2)$
and $(3\Lambda,7\Lambda/2)$ respectively, and those of the mirror family in the same intervals shifted by $4\Lambda$;
these sets are disjoint except for the last two. For different values of~$a$ the endpoints of $e_1$ differ, as do those
of $e_2$, of~$e_3$ and of~$e_4$, and $3\Lambda+a+2\eta=3\Lambda+a'$ exactly when $a'=a+2\eta$. So two members share a
hop only if their parameters differ by~$2\eta$, the shared hop being $e_3$ of the smaller and $e_4$ of the larger, and
within a member the four hops are distinct. Let $n=\Lambda/2-4\eta-1$ be the number of members of each family. The
members split into chains $a,a+2\eta,a+4\eta,\dots$; keeping every other member of each chain, starting with the
first, keeps $\lceil\ell/2\rceil$ members of a chain of~$\ell$, hence at least $n/2$ members of each family, and the
kept members are pairwise hop-disjoint. With the mirror family this gives at least~$n$ hop-disjoint cycles per period.

\emph{Counting.} On the $r$-fold cover, take these cycles in each of the $r$ period copies, translated by multiples of
$8\Lambda$; they remain hop-disjoint, each has defect $-4$ computed with~$\kappa_1$, and their number $n_{\mathrm c}$ is
at least $rn$. Interior cells have $4r$ visits per cover and $q\mid4r$, so by \cref{lem:chain}(c) the label~$G$
satisfies the hypothesis of \cref{thm:E}(b) with $\kappa=\kappa_1$. That theorem gives
$q\cost(G)\ge2T+n_{\mathrm c}(q-4)$, that is, $\cost(G)\ge n_{\mathrm c}+(2T-4n_{\mathrm c})/q\ge n_{\mathrm c}$,
because the $n_{\mathrm c}$ cycles use $4n_{\mathrm c}$ distinct hops among the $T$ hops of head~$2$. Since
$n_{\mathrm h}=2T=16r\Lambda$ and $n_{\mathrm c}\ge rn$, this is $\cost(G)\ge\nu_1n_{\mathrm h}$. If $q>4$, the right
side $n_{\mathrm c}(1-4/q)+2T/q$ increases with~$n_{\mathrm c}$, and $n_{\mathrm c}\ge2T\nu_1$ gives
$\cost(G)\ge\nu_1n_{\mathrm h}+(1-4\nu_1)n_{\mathrm h}/q$.

\emph{The golden shift.} Here $P=8\Lambda$ and $0.381966\,P=3.055728\Lambda$. For $s_0=5,6,7,8$ one finds
$\lfloor3.055728\Lambda\rfloor=97,195,391,782$, so $\xi=97,195,391,783$, $\eta=1,3,7,15$, $8\eta<\Lambda$, and
$\nu_1=11/512$, $19/1024$, $35/2048$, $67/4096$. For $s_0\ge9$, $3.055728\Lambda-1<\xi\le3.055728\Lambda+1$, so
$0.055728\Lambda-1<\eta\le0.055728\Lambda+1$; hence $0<8\eta<\Lambda$ and
$\nu_1\ge(0.277088-5/\Lambda)/16\ge0.0167>67/4096$, as $\Lambda\ge512$.
\end{informalproof}

\paragraph{Why it matters.} Theorem~E$_1$ is the first price per hop that does not depend on the depth, on the
sharpest open instance, with one level of nesting, for a natural class of labels: the class that extends the mechanism solving one walk by an
arbitrary correction per cell. On the golden pair with $m=1$ and $s_0\ge5$, such labels cost at least $0.016$ per hop on the cover, about
$0.033\,T$, at every depth~$D$, so they cannot give $\cost\le\eps T$ for $\eps<0.03$; for the aligned shift $\xi=1$ the
analogous certificates are close to zero numerically. The theorem says nothing about labels with small but nonzero
errors on
head~$1$, and these are exactly what optimised labels use (\cref{rem:dep-absorb}).

\begin{remark}[From the cover to the walk DAG]\label{rem:dep-linear}
The same argument works on the walk DAG of the pair of \cref{thm:E1} on $[0,NP)$, $N\ge1$, for labels
$\Psi\colon[0,NP)\to\{0,\dots,q-1\}$ that are exact modulo~$q$ on all head-$1$ hops at interior cells. The chains of
hops no longer close up, so summing $q[\Psi(u)\le\Psi(p)]=1+e-(\Psi(u)-\Psi(p))$ over the hops of a chain gives
$q\cdot\#\{\text{violated hops}\}\ge\#\{\text{hops}\}+\sum e-(q-1)\#\{\text{chains}\}$; there are at most
$2(2\Lambda+1)$ chains. Along each interior head-$1$ chain
$\Psi-\kappa_1$ is constant modulo~$q$, with no divisibility condition, and the cycles of the family that lie in
$[0,NP)$, at least $Nn$ after selection, force head-$2$ error at least $q-4$ each. Hence, counting hops per head,
$\cost(\Psi)\ge2\nu_1NP-O(\Lambda)$ for $q>4$, and $\cost(\Psi)\ge NP/2-O(\Lambda)$ for $q\le4$ from the floor. A
pair $(p,u)$ that is a hop of both heads would be a single edge of the walk DAG, but there is none. Indeed, all turn
times and turn positions of $U(s_0,m,g)$ are multiples of $2^{s_m}$, so on every maximal monotone run $x(t)=\pm t+c$
with $2^{s_m}\mid c$. A hop $(p,u)$ leaves $x(p)$ at time~$p$ and returns at time~$u$ along runs of opposite slopes,
so $p+u\equiv0\pmod{2^{s_m}}$, while a hop of $x_2$ satisfies $p+u\equiv-2\xi$; and $2\xi\not\equiv0\pmod{2^{s_m}}$ for
odd~$\xi$ when $s_m\ge2$; for the pair of \cref{thm:E1}, where $\xi=3\Lambda+\eta$ need not be odd, $2^{s_1}=\Lambda/2$
and $2\xi\equiv2\eta\not\equiv0\pmod{\Lambda/2}$ because $0<2\eta<\Lambda/4$.
\end{remark}

\begin{remark}[Numerical: nesting depths $m\le6$]\label{rem:dep-Em}
Theorem~E$_1$ is a closed form of computer-found certificates. For the golden pair $\mathrm{DNP}(s_0,m,1)$ with $1\le
m\le6$, explicit one-parameter families of negative cycles in~$M_2$, fitted as affine functions of~$a$, $\Lambda$
and~$\xi$ on ranges of~$s_0$ between $9$ and~$14$ and checked member by member on instances outside the fitting range
(up to $s_0=16$), certify for labels exact on all head-$1$ hops at interior cells, at every~$q$, prices per hop of
$0.0509$, $0.0554$, $0.0527$, $0.0519$, $0.0526$, $0.0480$ for $m=1,\dots,6$ (from $12,16,95,208,245,531$ families,
with the fractional packing that gives each member $C$ the weight $1/\max_{h\in C}\mathrm{load}(h)$, where the load of a
hop is the number of members through it). These values were stable in~$s_0$; they are per-instance certificates, not proofs
for all~$s_0$, and they have not been checked independently. For $m=1$ a greedy selection of hop-disjoint
members does better: it finds exactly $\Lambda-4$ negative $4$-cycles per period on all instances checked ($7\le s_0\le14$
and $s_0=16$), that is, a price $1/16-1/(4\Lambda)$ per hop at every~$q$; a proof would need the coincidences between the
twelve families. For labels exact only on
the head-$1$ hops of some levels, the certified price falls as levels are released: at $m=1$ exactness on the finest
level alone admits an explicit label of cost at most about $4/q$ per hop, while at $m=3$ and $s_0=7$, exactness on the
finest one, two or three levels or on all four carries certified prices of about $0.003$, $0.026$, $0.066$ and
$0.107$. Why ``every~$m$'' is open: the families show no recursion in~$m$, and the phases of the shift relative to the
level grids, $\psi_i=\mathrm{frac}(2^i\xi/2^{s_0})$, which are the same for all strokes of level~$i$, degenerate for
special~$m$; for instance $\psi_2,\dots,\psi_6$ lie between $0.003$ and $0.06$ at $m=8$, and when $15625$ divides
$m+1$ exact arithmetic gives $\psi_i=2^{i-s_0}$ for every $i\ge3$. A proof for all~$m$ therefore needs a rule
generating level-$i$ families from local configurations, a bound on their overlaps independent of~$m$, and a positive
share of good levels, not equidistribution of phases.
\end{remark}

\subsection{The remaining gap}\label{sec:dep-gap}

Labels found by optimisation are not exact on either head. They trade errors of head~$1$ on backward stretches
against errors of head~$2$ (\cref{rem:dep-absorb}), and by \cref{lem:transfer}(c) such an absorption costs only about
$\abs\delta/q$ per cycle. A lower bound for all labels must show that absorptions cannot be supplied consistently: an
error on a head-$1$ hop lies on backward stretches of some head-$2$ cycles, where it helps, and on forward stretches of
others, where it hurts, and symmetrically for head~$2$. We expect that this forces a multiplicative loss at every level
that both heads use, and we state the expected outcome as a count of served levels.

Let the pair of \cref{def:dnp} be defined, verbatim, for any odd shift~$\xi$. For a label~$G$ on its cover and
$0<\theta<1$, say that level~$i$ is \emph{$\theta$-served for head~$h$} if the hops of head~$h$ of level~$i$ have total error at most $\theta qn_{\mathrm h}/(2(m+1))$, and
\emph{doubly $\theta$-served} if it is $\theta$-served for both heads. When the $n_{\mathrm h}/2$ hops of each head are
spread evenly over the $m+1$ levels, as numerical counts indicate, this says that the mean error per hop of the class
is at most~$\theta q$. Two kinds of labels serve every level for trivial reasons; the conjecture excludes the first by
its bound on the cost and the second by its choice of the shift. Labels with few or random values serve every level at
large~$\theta$, but cost at least $n_{\mathrm h}/q$, more than $n_{\mathrm h}/8$ if $q<8$, respectively about
$n_{\mathrm h}/2$, by \eqref{eq:dep-Zq}. And for a shift aligned with the dyadic scales, such as the golden shift
when $15625$ divides $m+1$, a label exact on one head is nearly exact on the other.

\begin{conjecture}[Conjecture A$'$: few doubly served levels]\label{conj:Aprime}
There are $\theta_0>0$ and constants $C,C'$, and for all $s_0$ and~$m$ an odd shift~$\xi$, such that for all $r$
and~$q$ every label~$G$ with $\cost(G)\le n_{\mathrm h}/8$ on the $r$-fold cover of the pair $x_1=U(s_0,m,1)$,
$x_2=x_1(\cdot+\xi)$ has at most $C\log_2q+C'$ doubly $\theta_0$-served levels.
\end{conjecture}

We expect shifts misaligned with the dyadic scales at every level to qualify. The golden shift is not always one of
them: even when $15625\nmid m+1$, its phases nearly align at several levels for some~$m$, such as $m=8$
(\cref{rem:dep-Em}).

If the conjecture holds, let $G$ be a label, on a cover of the pair with the shift that it provides, of cost
$\cost(G)\le\eps n_{\mathrm h}$, where $\eps\le1/8$ and
$\eps<\theta_0/2$. By \cref{cor:dep-split}, every level that is not doubly $\theta_0$-served contributes more than
$\theta_0n_{\mathrm h}/(2(m+1))$ to the cost, so fewer than $2\eps(m+1)/\theta_0$ levels are not doubly served, and
$(m+1)(1-2\eps/\theta_0)<C\log_2q+C'$. This gives
\[
  D=q-1\;\ge\;\exp\bigl(c(1-2\eps/\theta_0)(m+1)-c'\bigr)-1
\]
for constants $c,c'>0$, so the depth needed at a fixed cost per hop grows exponentially in the depth of nesting,
matching the clock labels of \cref{prop:dep-clock} up to the constant in the exponent. The same statement for labels
on the walk DAG of many periods would bound the depth of every cut of size $\eps T$ in the same way, because
\eqref{eq:dep-Zq} and \cref{cor:dep-split} hold there up to boundary terms and, when $s_m\ge2$, no pair is a hop of
both heads (both as in \cref{rem:dep-linear}). It would make the family of these pairs with $s_m\ge2$, $m\to\infty$,
depth-robust, and \cref{prob:tcp} for $k=2$ would have a negative answer.

\begin{remark}[Numerical: evidence for Conjecture~A$'$]\label{rem:dep-Aevidence}
The evidence is limited. In the best labels found by optimisation for $m=3$ ($g=1$, $s_0=6$), with mean error at most
$q/10$ per hop counted as served, the number of doubly served levels was $1$, $1$, $2$ and $2$ or~$3$ at
$D=7,15,31,63$, consistent with growth like $\log q$, and level~$0$ of head~$1$ was never served (mean error about
$0.17q$ to $0.20q$ at every~$D$). Labels exact on several levels of one head carry a price that grows with the number
of exact levels (\cref{rem:dep-Em}), a restricted form of the conjecture. Small instances cannot decide asymptotic
questions, and these labels are upper bounds of uncertain quality.
\end{remark}

\paragraph{Status.} We do not know whether the deep-nesting pairs are depth-robust as $m$ grows. No mechanism we know
cuts them at a depth independent of~$m$: per-head band rules appear to fail (proof sketches of \cref{rem:dep-N,rem:dep-Nprime}); for $g\ge2$, anchored digits
cannot be monotone (\cref{lem:dep-monotone}); and explicit combinations of visit indices and band counts kept a cost per
hop bounded away from zero at every depth tried, in computations up to $m=8$. Clock labels cut each pair, but only at depth $(c/\eps)^{m+1}$ (\cref{prop:dep-clock}). On the other
side, the only lower bounds that hold at every depth are for restricted classes
(\cref{thm:E,thm:E1}), and packing certificates cannot do better (\cref{prop:packing}).

\paragraph{Why it matters.} \Cref{prob:tcp} is the walk-level form of the structural hypothesis of \cref{thm:D}, and
\cref{thm:B} solves it for home-based walks. On the sharpest open instance, the evidence just listed leaves a robust
multiplicative count of absorptions as the missing step, of which \cref{conj:Aprime} is one form. Either answer is
informative. If the pairs are depth-robust as $m$ grows, $\TCP(2)$ fails as stated; if moreover oblivious linear-time
machines can generate them as dummy traffic and they are depth-robust at every $\eps$ below the crux rates in question, the structural hypothesis of \cref{thm:D} must be weakened as explained in
\cref{sec:dep-problem}, and these pairs, placed at intermediate scales inside curved motion, are candidates for the depth-robust schedule that
an oblivious linear-time transposer obeying the crux would need. If they are depth-reducible at a depth independent
of~$m$, the mechanism must beat clock
labels by absorbing defects consistently across levels, and it would be the core of a general proof of $\TCP(2)$.
